% archon:covers AlgebraicJacobian/Albanese/CodimOneExtension.lean
% NOTE: Weil-divisor API extension of RiemannRoch_WeilDivisor.tex --- re-uses
% Scheme.WeilDivisor, Scheme.PrimeDivisor, Scheme.RationalMap.order from
% \cref{chap:RiemannRoch_WeilDivisor}. The order map
% \texttt{AlgebraicGeometry.Scheme.RationalMap.order} is not re-defined here;
% see \cref{def:order_at_point} for the canonical reference.

\chapter{Codimension-1 indeterminacy extension (A.4.a)}
\label{chap:Albanese_CodimOneExtension}

% STRATEGY NOTE (iter-174). This chapter is the A.4.a sub-row of Route A
% (the positive-genus arm of \texttt{nonempty\_jacobianWitness}). Per
% STRATEGY.md L29 it is the \emph{risk-dominant} A.4 sub-row
% ("project-fatal if it stalls"). The chapter has two outputs:
% (i) \cref{thm:codim_one_extension} --- Milne Theorem~3.1: a rational map
%     from a normal variety to a complete variety is defined on the complement
%     of a codimension-\(\geq 2\) closed subset. Specialised below to the
%     smooth-variety setting where the prover can phrase it as
%     "no codim-1 indeterminacy \(\Rightarrow\) extension". This is one of the
%     two inputs to Milne Theorem~3.2 in the sibling chapter
%     \cref{chap:Albanese_Thm32RationalMapExtension}.
% (ii) \cref{lem:milne_codim1_indeterminacy} --- Milne Lemma~3.3: for a
%     rational map from a nonsingular variety to a group variety, the
%     indeterminacy locus is either empty or of pure codimension \(1\). This is
%     the second input to Milne Theorem~3.2; it consumes nothing from A.4.b.
% A bonus output is \texttt{thm:weil_divisor_obstruction}: the Weil-divisor
% reformulation of (i) that a smooth-surface codim-2 extension is obstructed
% exactly by the order of the pulled-back coordinates of \(f\) being negative
% at a prime divisor. This is the lemma A.4.a uses to align with the project's
% \cref{chap:RiemannRoch_WeilDivisor} machinery.
%
% Depth-\(\geq 2\) extension on a smooth surface, used in the proof body of
% \cref{thm:codim_one_extension} when reducing the general-codim case to the
% normal-DVR-codim-1 case, consumes
% \cref{cor:regular_cohen_macaulay} of
% \cref{chap:Albanese_AuslanderBuchsbaum} (the A.4.b sub-row). The chapter
% records the dependency but does not re-prove it.
%
% The Lean target file \texttt{AlgebraicJacobian/Albanese/CodimOneExtension.lean}
% does not yet exist; this chapter is the prover-ready specification for
% creating it.

\section{Setup and motivation}
\label{sec:codim1_setup}

Fix an algebraically closed field \(\bar k\) throughout. By a \emph{variety} over \(\bar k\)
we mean a geometrically integral, separated scheme of finite type over \(\bar k\)
(Milne's convention, \emph{Abelian Varieties}, p.~v); a variety is \emph{nonsingular}
(or \emph{smooth}) if its local ring at every closed point is regular. A variety is
\emph{complete} if it is proper over \(\bar k\).

A \emph{rational map} \(f \colon X \dashrightarrow Y\) between varieties is the equivalence
class of a regular morphism \(\varphi_U \colon U \to Y\) defined on a dense open subset
\(U \subseteq X\), where two pairs \((U, \varphi_U)\), \((U', \varphi_{U'})\) are equivalent
if \(\varphi_U\) and \(\varphi_{U'}\) agree on \(U \cap U'\) (Milne~\S I.3, p.~16). The map \(f\)
is \emph{defined at} a point \(x \in X\) if there is some representative \((U, \varphi_U)\)
with \(x \in U\); the set of such \(x\) is an open subset \(\mathrm{Dom}(f) \subseteq X\),
the \emph{domain of definition} of \(f\).

The motivating question of this chapter, following Milne~\S I.3, is:
\begin{quote}
  \emph{When does a rational map \(f \colon X \dashrightarrow Y\) extend to a regular
  morphism on all of \(X\)?}
\end{quote}
Without hypotheses on \(X\) or \(Y\) the question has no positive answer (Milne lists three
counterexamples on p.~16: the inclusion of a proper open into the ambient variety; the
rational map to the cuspidal cubic whose inverse fails to extend; the inverse of a
blow-up at a smooth point). The two structural inputs that buy unconditional extension
in our setting are:
\begin{itemize}
  \item \emph{Target completeness:} if \(Y\) is complete (proper over \(\bar k\)), the
    valuative criterion of properness forces the rational map to extend across every
    codim-\(1\) point of a normal source \(X\) (Milne Theorem~3.1).
  \item \emph{Target group structure:} if \(Y\) is a group variety, the difference map
    \(\Phi(x, y) := f(x) \cdot f(y)^{-1}\) on \(X \times X\) has a vanishing/pole divisor
    along the diagonal that controls the codimension of the indeterminacy locus
    (Milne Lemma~3.3).
\end{itemize}
For \(Y\) an abelian variety both conditions hold, and combining them yields Milne's
Theorem~3.2 (the headline of \cref{chap:Albanese_Thm32RationalMapExtension}). The
present chapter packages the two inputs separately:
\cref{thm:codim_one_extension} is Milne~3.1, and
\cref{lem:milne_codim1_indeterminacy} is Milne~3.3. The Weil-divisor reformulation
\texttt{thm:weil_divisor_obstruction} translates the codim-\(1\) extension obstruction
in Milne~3.1's proof into the order-at-a-prime-divisor language pinned in
\cref{chap:RiemannRoch_WeilDivisor}, which is the form the Lean prover consumes when
working with the project's \(\mathrm{Scheme.WeilDivisor}\) API.

\section{The indeterminacy locus of a rational map}
\label{sec:indeterm_locus}

\begin{definition}


\leanok
  [Indeterminacy locus of a rational map]
  \label{def:indeterminacy_locus}
  \lean{AlgebraicGeometry.Scheme.RationalMap.indeterminacyLocus}
  % SOURCE: [Milne, Abelian Varieties], \S I.3, p.~16, the paragraph "Let \(V\) and \(W\)
  % be varieties over \(k\)\ldots" defining the domain of definition of a rational map
  % (read from references/abelian-varieties.pdf, PDF page 22).
  % SOURCE QUOTE: "A rational map \(\varphi\) is said to be \emph{defined} at a point
  % \(v\) of \(V\) if \(v \in U\) for some \((U, \varphi_U) \in \varphi\). The set \(U_1\) of
  % \(v\) at which \(\varphi\) is defined is open, and there is a regular map
  % \(\varphi_1 \colon U_1 \to W\) such that \((U_1, \varphi_1) \in \varphi\) --- clearly,
  % \(U_1 = \bigcup_{(U, \varphi_U) \in \varphi} U\) and we can define \(\varphi_1\) to be
  % the regular map such that \(\varphi_1 | U = \varphi_U\) for all $(U, \varphi_U) \in
  % \varphi$."
  \textit{Source: Milne, \emph{Abelian Varieties}, \S I.3, p.~16.}
  Let \(f \colon X \dashrightarrow Y\) be a rational map between varieties over \(\bar k\).
  The \emph{domain of definition} of \(f\) is the open set
  \[
    \mathrm{Dom}(f) \;:=\; \bigcup_{(U, \varphi_U) \in f} U \;\subseteq\; X,
  \]
  i.e.\ the union of the dense opens of all representatives of \(f\). The
  \emph{indeterminacy locus} of \(f\) is the closed complement
  \[
    Z(f) \;:=\; X \smallsetminus \mathrm{Dom}(f) \;\subseteq\; X.
  \]
  The map \(f\) is regular (i.e.\ defined on all of \(X\)) if and only if \(Z(f) = \emptyset\).

  \paragraph{Lean encoding scope.} The Lean target
  \texttt{AlgebraicGeometry.Scheme.RationalMap.indeterminacyLocus} is a method on the
  Mathlib \texttt{Scheme.RationalMap} structure pinned at
  \texttt{Mathlib.AlgebraicGeometry.Birational.RationalMap}; it returns a
  \texttt{Set X} (the underlying carrier of the closed complement of
  \texttt{Scheme.RationalMap.domain}, which is already exposed by Mathlib). The
  closedness of \(Z(f)\) is a separate \texttt{IsClosed} witness, packaged either as a
  field of \texttt{indeterminacyLocus} or as a standalone lemma
  \texttt{Scheme.RationalMap.isClosed\_indeterminacyLocus}.
\end{definition}


\begin{definition}


\leanok
  [Codimension-1-indeterminacy-free rational maps]
  \label{def:codim_one_indeterminacy}
  \lean{AlgebraicGeometry.Scheme.RationalMap.CodimOneFree}
  \uses{def:indeterminacy_locus, def:prime_divisor}
  % SOURCE: [Milne, Abelian Varieties], Theorem~3.1, \S I.3, p.~16 --- the conclusion
  % "an open subset \(U\) of \(V\) whose complement \(V \smallsetminus U\) has codimension \(\geq 2\)"
  % is the standing predicate we abstract (read from references/abelian-varieties.pdf,
  % PDF page 22).
  % SOURCE QUOTE: "A rational map \(\varphi \colon V \dashrightarrow W\) from a normal
  % variety \(V\) to a complete variety \(W\) is defined on an open subset \(U\) of \(V\) whose
  % complement \(V \smallsetminus U\) has codimension \(\geq 2\)."
  \textit{Source: Milne, \emph{Abelian Varieties}, Theorem~3.1, \S I.3, p.~16 (the
  conclusion).}
  Let \(f \colon X \dashrightarrow Y\) be a rational map between varieties over \(\bar k\).
  We say \(f\) is \emph{codim-1-indeterminacy-free} (or has \emph{codimension-\(\geq 2\)
  indeterminacy}) if every irreducible component of the closed subset \(Z(f) \subseteq X\)
  has codimension \(\geq 2\) in \(X\). Equivalently, no prime divisor of \(X\) is contained in
  \(Z(f)\); equivalently, every point \(\eta \in X\) with \(\mathrm{codim}_X \overline{\{\eta\}} = 1\)
  lies in \(\mathrm{Dom}(f)\).

  \paragraph{Why the abstraction.} Milne's Theorem~3.1 (\cref{thm:codim_one_extension}
  below) and Lemma~3.3 (\cref{lem:milne_codim1_indeterminacy} below) both produce
  conclusions about the codimension of \(Z(f)\). Phrasing the codim-\(\geq 2\) condition
  as a named predicate keeps the cited statements concise and lets the consumer
  \cref{chap:Albanese_Thm32RationalMapExtension} read off the combination
  (\(Z(f)\) has codim \(\geq 2\) from \cref{thm:codim_one_extension} \emph{and} \(Z(f)\) is
  empty or pure codim 1 from \cref{lem:milne_codim1_indeterminacy} \(\Rightarrow\) \(Z(f)\)
  is empty) as a single line.

  \paragraph{Lean encoding scope.} The Lean target
  \texttt{AlgebraicGeometry.Scheme.RationalMap.CodimOneFree} is a \texttt{Prop}-valued
  predicate, encoded via Mathlib's \texttt{Order.coheight} on the specialisation
  preorder of \texttt{X.carrier} as in
  \cref{def:prime_divisor}: \(f\) is codim-1-indeterminacy-free if and only if for every
  \(\eta \in X\) with \(\texttt{Order.coheight\ }\eta = 1\), \(\eta \in
  \texttt{f.domain}\). The encoding matches the standing \(\mathrm{Order.coheight}\)
  convention pinned in \cref{def:prime_divisor}, so prime divisors of \(X\) and the
  hypothesis ``no codim-\(1\) indeterminacy'' use the same Mathlib idiom.
\end{definition}

\section{Smoothness yields a DVR at every codim-1 point}
\label{sec:smooth_dvr}

The codim-\(1\) extension argument in Milne's proof of Theorem~3.1 hinges on the local
ring at the generic point of a prime divisor being a discrete valuation ring (so that
the valuative criterion of properness applies). On a normal variety this is the
defining property of regularity in codimension one; on a smooth variety it is automatic.

The smooth-implies-regular-local-stalk implication itself is the Jacobian-criterion
direction of Stacks tag \texttt{00TT}; we package it as a separate sub-lemma so that
\cref{lem:smooth_codim_one_dvr} below reads cleanly as ``regular local of dimension
\(1\) is a DVR'', with the regularity prerequisite supplied by the sub-lemma.

\subsection{K\"ahler-differentials localisation substrate}
\label{subsec:kaehler_localisation_substrate}

The Jacobian-criterion direction of Stacks tag \texttt{00TT}
(\cref{lem:smooth_to_regular_local_ring} below) is assembled from several
substrate results: smoothness gives flatness at a stalk and, locally, a
standard-smooth presentation; the scheme-to-algebra bridge and the freeness of
\(\Omega_{S/R}\) for a standard-smooth algebra \(R \to S\) then reduce the
regularity claim to a freeness-plus-rank computation on K\"ahler differentials,
transported through localisation at the prime defining the stalk. The two
localisation helpers below --- transport of freeness, and matching of rank under
localisation --- supply this reduction. The Krull-dimension and cotangent inputs
are collected in \cref{subsec:stage6_subgap_decomposition}, and the regularity
conclusion at an arbitrary prime is \cref{thm:standard_smooth_regular_at_prime}.

The substrate lemmas are recorded below, feeding
\cref{lem:smooth_to_regular_local_ring}.

\begin{lemma}\leanok
[Stage 1 --- smooth implies flat at every stalk]
  \label{lem:stalkMap_flat_of_smooth}
  \lean{AlgebraicGeometry.Scheme.stalkMap_flat_of_smooth}
  Let \(X \to \Spec(\bar k)\) be a smooth morphism over an algebraically closed
  field \(\bar k\). Then for every point \(z \in X\) the induced ring map on
  stalks is flat.
\end{lemma}

\begin{proof}
  Proved directly in Lean: the smooth-implies-flat instance composed with the
  stalk-map flatness characterisation. A sub-step of
  \cref{lem:smooth_to_regular_local_ring}.
\end{proof}

\begin{lemma}\leanok
[Stage 2 --- smooth implies a standard-smooth presentation at a point]
  \label{lem:exists_standardSmooth_at_of_smooth}
  \lean{AlgebraicGeometry.Scheme.exists_isStandardSmooth_at_of_smooth}
  For a smooth morphism \(X \to \Spec(\bar k)\) and a point \(z \in X\), there
  exist affine open neighbourhoods \(U \subseteq \Spec(\bar k)\) and \(V \ni z\)
  in \(X\) with \(V \subseteq X^{-1}(U)\) such that the induced ring map
  \(\Gamma(\Spec \bar k, U) \to \Gamma(X, V)\) is standard smooth (Stacks 00T7).
\end{lemma}

\begin{proof}
  Proved directly in Lean: re-export of the scheme-level Stacks 00T7 packaging.
  A sub-step of \cref{lem:exists_algebra_standardSmooth_stalk_localization}.
\end{proof}

\begin{lemma}\leanok
[Stage 3 --- standard-smooth algebra package at the stalk]
  \label{lem:exists_algebra_standardSmooth_stalk_localization}
  \lean{AlgebraicGeometry.Scheme.exists_algebra_isStandardSmooth_section_stalk_isLocalization_of_smooth}
  \uses{lem:exists_standardSmooth_at_of_smooth, lem:isLocalization_atPrime_stalk_of_affineOpen}
  For a smooth morphism \(X \to \Spec(\bar k)\) and a point \(z \in X\), there is
  an affine open neighbourhood \(V \ni z\) and an \(\Gamma(\Spec \bar k, U)\)-algebra
  structure on \(\Gamma(X, V)\) making it a standard-smooth algebra, together with a
  witness that the stalk \(\mathcal O_{X, z}\) is the localisation of \(\Gamma(X, V)\)
  at the prime ideal corresponding to \(z\).
\end{lemma}

\begin{proof}
  Proved directly in Lean: composes the Stage-2 presentation
  \cref{lem:exists_standardSmooth_at_of_smooth} with the ring-hom-to-algebra bridge
  and the affine-open stalk localisation
  \cref{lem:isLocalization_atPrime_stalk_of_affineOpen}. A sub-step of
  \cref{lem:smooth_to_regular_local_ring}.
\end{proof}

\begin{lemma}\leanok
[Stage 4 --- K\"ahler differentials of a standard-smooth algebra are free]
  \label{lem:module_free_kaehler_standardSmooth}
  \lean{AlgebraicGeometry.Scheme.module_free_kaehlerDifferential_of_isStandardSmooth}
  For a standard-smooth \(R\)-algebra \(S\), the module of K\"ahler differentials
  \(\Omega_{S/R}\) is free as an \(S\)-module (Stacks 00T7(2)), with basis the family
  \(\{ds_i\}\) of the standard-smooth presentation.
\end{lemma}

\begin{proof}
  Proved directly in Lean: from the basis provided by the standard-smooth
  presentation. Supplies the freeness hypothesis of
  \cref{lem:module_free_kaehler_localization}.
\end{proof}

\begin{lemma}


\leanok
  [Stage 5a --- K\"ahler-differentials freeness transports through localisation]
  \label{lem:module_free_kaehler_localization}
  \lean{AlgebraicGeometry.Scheme.module_free_kaehlerDifferential_localization}
  \uses{lem:module_free_kaehler_standardSmooth}
  Let \(R\) and \(S\) be commutative rings equipped with an \(R\)-algebra
  structure on \(S\), and assume that the module of K\"ahler differentials
  \(\Omega_{S/R}\) is free as an \(S\)-module. Let \(M \subseteq S\) be a
  submonoid and let \(S_M\) be a localisation of \(S\) at \(M\), so that
  \(R \to S \to S_M\) is a tower of \(R\)-algebras. Then the module of
  K\"ahler differentials \(\Omega_{S_M / R}\) is free as an
  \(S_M\)-module.
\end{lemma}

\begin{proof}
  The canonical comparison map
  \(\Omega_{S/R} \to \Omega_{S_M/R}\) induced by the localisation map
  \(S \to S_M\) realises \(\Omega_{S_M/R}\) as the localisation of the
  \(S\)-module \(\Omega_{S/R}\) at \(M\); i.e.\ it satisfies Mathlib's
  \texttt{IsLocalizedModule} predicate along the underlying ring
  localisation. This is Stacks tag \texttt{02JK}, packaged in Mathlib as
  \texttt{KaehlerDifferential.isLocalizedModule\_map}. Freeness of a
  module is preserved under localisation of modules via Mathlib's
  \texttt{Module.free\_of\_isLocalizedModule}; combining the two yields
  the conclusion.
\end{proof}

\begin{lemma}


\leanok
  [Stage 5b --- K\"ahler-differentials rank under localisation equals
  the relative dimension]
  \label{lem:rank_kaehler_localization_eq_relative_dim}
  \lean{AlgebraicGeometry.Scheme.rank_kaehlerDifferential_localization_eq_relativeDimension}
  \uses{lem:module_free_kaehler_localization}
  With notation as in \cref{lem:module_free_kaehler_localization},
  assume in addition that the \(R\)-algebra \(S\) is \emph{standard
  smooth of relative dimension \(n\)} (Mathlib's
  \texttt{Algebra.IsStandardSmoothOfRelativeDimension}\,\(n\)\,\(R\)\,\(S\),
  the algebra-side packaging of Stacks tag \texttt{00T7}(2)) and that
  the localisation \(S_M\) is non-trivial as a ring. Then the rank of
  \(\Omega_{S_M / R}\) over \(S_M\) equals \(n\).
\end{lemma}

\begin{proof}
  Mathlib's
  \texttt{Algebra.IsStandardSmoothOfRelativeDimension.\linebreak rank\_kaehlerDifferential}
  asserts \(\mathrm{rank}_S(\Omega_{S/R}) = n\) under the
  standard-smooth hypothesis. The localisation map
  \(\Omega_{S/R} \to \Omega_{S_M/R}\) is an
  \texttt{IsLocalizedModule} along \(S \to S_M\) (Stacks tag
  \texttt{02JK}; cf.\ the proof of
  \cref{lem:module_free_kaehler_localization}), and Mathlib's
  \texttt{Module.lift\_rank\_of\_isLocalizedModule\_of\_free} transports
  module rank through an \texttt{IsLocalizedModule} up to
  \(\mathrm{Cardinal.lift}\). Since \(n\) is a natural number both
  ranks equal \(n\), so
  \(\mathrm{rank}_{S_M}(\Omega_{S_M / R}) = n\) as required.
\end{proof}

\subsection{Stage 6 sub-decomposition: 00OE + 02JK chain}
\label{subsec:stage6_subgap_decomposition}

Given a standard-smooth presentation of an affine chart \(R \to S\) at a
point \(z\), with \(\Omega_{S/R}\) free of rank \(n\) and
\(\Omega_{S_\mathfrak q / R}\) free of rank \(n\) on the stalk
\(S_\mathfrak q = \mathcal O_{X, z}\), one concludes that \(S_\mathfrak q\)
is a regular local ring from two ingredients: (6.A) the smooth-algebra
Krull-dimension formula (Stacks tag \texttt{00OE}), and (6.B) the cotangent /
K\"ahler-differential bridge over a field base (Stacks tag \texttt{02JK}). The
two ingredients combine (6.C) into the regularity conclusion via
\(\dim S_\mathfrak q = \mathrm{rank}_{S_\mathfrak q}\Omega_{S_\mathfrak
q / R} = \dim_{\kappa(\mathfrak q)} \mathfrak m / \mathfrak m^2\),
the cotangent-space characterisation of regularity. The next three lemmas
record these three ingredients.

\begin{lemma}
  [Stage 6.A --- Smooth-algebra Krull-dimension formula (Stacks 00OE)]
  \label{lem:smooth_algebra_krull_dim_formula}
  \uses{lem:rank_kaehler_localization_eq_relative_dim,
        lem:mvPolynomial_height_eq_natCard_of_isMaximal,
        lem:standard_smooth_le_height_of_isMaximal,
        lem:standard_smooth_le_ringKrullDim_localization,
        lem:matsumura_isRegular_of_linearIndependent_cotangent,
        lem:exists_submersivePresentation_of_standardSmoothReldim}
  % SOURCE: [Stacks Project], lemma-dimension-at-a-point-finite-type-field
  % (canonically known as tag 00OE; tag-to-label cross-check not in the
  % local `tags/tags` file, see iter-198 writer report note) (read from
  % references/stacks-algebra.tex, L28206--28217 in section
  % "Dimension theory of finitely generated $k$-algebras").
  % SOURCE QUOTE: "Let $k$ be a field. Let $S$ be a finite type $k$
  % algebra. Let $X = \Spec(S)$. Let $\mathfrak p \subset S$ be a prime
  % ideal, and let $x \in X$ be the corresponding point. Then we have
  % $\dim_x(X) = \dim(S_{\mathfrak p}) + \text{trdeg}_k\
  % \kappa(\mathfrak p)$."
  \textit{Source: [Stacks Project], tag~00OE (\(=\) \texttt{lemma-dimension-at-a-point-finite-type-field}).}
  Let \(k\) be a field, let \(S\) be a finite-type \(k\)-algebra, and let
  \(\mathfrak q \subset S\) be a prime ideal corresponding to a point
  \(x \in \mathrm{Spec}(S)\). Then
  \[
    \dim_x(\mathrm{Spec}\,S)
    \;=\; \dim S_{\mathfrak q} + \mathrm{trdeg}_k\,\kappa(\mathfrak q).
  \]
  In particular, if \(\mathfrak q\) is a maximal ideal (so \(x\) is a
  closed point) and \(k\) is algebraically closed, then
  \(\kappa(\mathfrak q) = k\) (Nullstellensatz), so
  \(\mathrm{trdeg}_k\,\kappa(\mathfrak q) = 0\) and
  \(\dim S_{\mathfrak q} = \dim_x(\mathrm{Spec}\,S)\).

  \paragraph{Geometric application.} Applied to the standard-smooth
  setup --- \(k = \bar k\), \(S = \Gamma(X.\mathrm{left}, V)\) a
  standard-smooth \(\bar k\)-algebra of relative dimension \(n\),
  and \(\mathfrak q\) the prime defining the point \(z \in V\) --- the
  formula yields \(\dim(\mathcal O_{X, z}) = \dim S_{\mathfrak q} = n -
  \mathrm{trdeg}_{\bar k}\,\kappa(z)\). For \(z\) a closed point this is
  simply \(n\); for the generic point of a codim-\(1\) subvariety
  (\(\mathrm{trdeg} = n - 1\)) it is \(1\).
\end{lemma}

% SOURCE QUOTE PROOF: "By Lemma \texttt{lemma-dimension-prime-polynomial-ring}
% we know that $r = \text{trdeg}_k\ \kappa(\mathfrak p)$ is equal to the
% dimension of $V(\mathfrak p)$. Pick any maximal chain of primes
% $\mathfrak p \subset \mathfrak p_1 \subset \ldots \subset \mathfrak p_r$
% starting with $\mathfrak p$ in $S$. This has length $r$ by Lemma
% \texttt{lemma-dimension-spell-it-out}. Let $\mathfrak q_j$, $j \in J$ be
% the minimal primes of $S$ which are contained in $\mathfrak p$. These
% correspond $1-1$ to minimal primes in $S_{\mathfrak p}$ via the rule
% $\mathfrak q_j \mapsto \mathfrak q_jS_{\mathfrak p}$. By Lemma
% \texttt{lemma-dimension-at-a-point-finite-type-over-field} we know that
% $\dim_x(X)$ is equal to the maximum of the dimensions of the rings
% $S/\mathfrak q_j$. For each $j$ pick a maximal chain of primes
% $\mathfrak q_j \subset \mathfrak p'_1 \subset \ldots \subset
% \mathfrak p'_{s(j)} = \mathfrak p$. Then $\dim(S_{\mathfrak p}) =
% \max_{j \in J} s(j)$. Now, each chain
% $\mathfrak q_j \subset \mathfrak p'_1 \subset \ldots \subset
% \mathfrak p'_{s(j)} = \mathfrak p \subset \mathfrak p_1 \subset \ldots
% \subset \mathfrak p_r$ is a maximal chain in $S/\mathfrak q_j$, and by
% what was said before we have $\dim_x(X) = \max_{j \in J} r + s(j)$.
% The lemma follows."
\begin{proof}
  \uses{lem:rank_kaehler_localization_eq_relative_dim,
        lem:mvPolynomial_height_eq_natCard_of_isMaximal,
        lem:standard_smooth_le_height_of_isMaximal,
        lem:standard_smooth_le_ringKrullDim_localization,
        lem:matsumura_isRegular_of_linearIndependent_cotangent,
        lem:exists_submersivePresentation_of_standardSmoothReldim}
  This is exactly Stacks tag~\texttt{00OE}
  (\texttt{lemma-dimension-at-a-point-finite-type-field}); the proof is a
  citation of the dimension theory of finitely generated algebras over a
  field. By Noether normalisation the finite-type \(k\)-algebra \(S\) is
  finite over a polynomial subalgebra, so the going-up / catenary
  dimension theory of such algebras applies. Write
  \(r = \mathrm{trdeg}_k\,\kappa(\mathfrak q)\); this transcendence
  degree is exactly the dimension of the closed subvariety
  \(\overline{\{x\}} = V(\mathfrak q)\). A maximal chain of primes
  through \(\mathfrak q\) therefore splits into a
  length-\(\dim S_{\mathfrak q}\) part below \(\mathfrak q\) (recording
  the local codimension, i.e.\ the height of \(\mathfrak q\) in the
  relevant irreducible component) and a length-\(r\) part above it
  (recording the dimension of \(\overline{\{x\}}\)). Taking the maximum
  over the minimal primes contained in \(\mathfrak q\) yields
  \[
    \dim_x(\mathrm{Spec}\,S) = \dim S_{\mathfrak q} + r
    = \dim S_{\mathfrak q} + \mathrm{trdeg}_k\,\kappa(\mathfrak q),
  \]
  which is the asserted formula.

  For the geometric specialisation matching the Lean target, take
  \(k = \bar k\) algebraically closed and \(S\) standard smooth of
  relative dimension \(n\), so that \(\dim_x(\mathrm{Spec}\,S) = n\) via
  \cref{lem:rank_kaehler_localization_eq_relative_dim} (the free
  K\"ahler module \(\Omega_{S_{\mathfrak q}/R}\) has rank \(n\), which
  pins the relative dimension). The formula then gives
  \(\dim S_{\mathfrak q} = n - \mathrm{trdeg}_{\bar k}\,\kappa(\mathfrak
  q)\). At a closed point \(\mathfrak q\) is maximal, so by the
  Nullstellensatz \(\kappa(\mathfrak q) = \bar k\) and
  \(\mathrm{trdeg}_{\bar k}\,\kappa(\mathfrak q) = 0\); hence
  \(\dim S_{\mathfrak q} = \dim_x(\mathrm{Spec}\,S) = n\).
\end{proof}

\begin{lemma}
  [Stage 6.B --- Cotangent \(\leftrightarrow\) K\"ahler over a field (Stacks 02JK)]
  \label{lem:cotangent_kahler_over_field}
  \lean{AlgebraicGeometry.Scheme.cotangent_iso_residue_tensor_kaehler_of_formallySmooth_residue}
  \uses{lem:smooth_algebra_krull_dim_formula}
  % SOURCE: [Stacks Project], lemma-differential-seq (canonically
  % known as the right-exact conormal portion of tag 02JK; the
  % short-exact / split refinement is the proof step inside
  % lemma-separable-smooth) (read from
  % references/stacks-algebra.tex, L34087--34104 and L38724--38766).
  % SOURCE QUOTE: "In diagram (\texttt{equation-functorial-omega}),
  % suppose that $S \to S'$ is surjective with kernel $I \subset S$, and
  % assume that $R' = R$. Then there is a canonical exact sequence of
  % $S'$-modules
  % $$I/I^2 \longrightarrow \Omega_{S/R} \otimes_S S' \longrightarrow
  % \Omega_{S'/R} \longrightarrow 0$$
  % The leftmost map is characterized by the rule that $f \in I$ maps
  % to $\text{d}f \otimes 1$."
  % SOURCE QUOTE: "Hence we get
  % $$\dim_{\kappa} \Omega_{R/k} \otimes_R \kappa = \dim_\kappa
  % \mathfrak m/\mathfrak m^2 + \text{trdeg}_k (\kappa) \geq \dim R +
  % \text{trdeg}_k (\kappa) = \dim_{\mathfrak q} S$$
  % (see Lemma \texttt{lemma-dimension-at-a-point-finite-type-field} for
  % the last equality) with equality if and only if $R$ is regular."
  \textit{Source: [Stacks Project], tag~02JK (right-exact half = \texttt{lemma-differential-seq}; separable-residue refinement inside \texttt{lemma-separable-smooth}).}
  Let \(k\) be a field and let \(R\) be a Noetherian local
  \(k\)-algebra with maximal ideal \(\mathfrak m\) and residue field
  \(\kappa = R/\mathfrak m\). The canonical exact sequence
  \[
    \mathfrak m / \mathfrak m^2
    \;\longrightarrow\; \Omega_{R/k} \otimes_R \kappa
    \;\longrightarrow\; \Omega_{\kappa / k}
    \;\longrightarrow\; 0
  \]
  identifies the cotangent space \(\mathfrak m / \mathfrak m^2\) (a
  \(\kappa\)-vector space) with the kernel of the right map. When
  \(\kappa / k\) is separable (in particular when \(k\) is
  algebraically closed and \(\kappa = k\), so \(\Omega_{\kappa / k}
  = 0\)) the leftmost map is injective and the sequence is short exact
  --- in fact split short exact when there is a \(k\)-algebra section
  \(\kappa \to R\) (e.g.\ at a closed point of a \(\bar k\)-variety).
  In particular, over an algebraically closed field \(k = \bar k\) and
  at a closed point \(\mathfrak m\),
  \[
    \Omega_{R/k} \otimes_R \kappa \;\simeq\; \mathfrak m / \mathfrak m^2
  \]
  as \(\kappa\)-vector spaces; hence
  \(\dim_\kappa(\Omega_{R/k} \otimes_R \kappa)
   = \dim_\kappa \mathfrak m / \mathfrak m^2\).
\end{lemma}

The Stage 6.B cotangent--K\"ahler chain is realised on the Lean side by the
following substrate lemmas, which compute the residue-field tensor finrank and
promote the cotangent isomorphism of \cref{lem:cotangent_kahler_over_field} into
the cotangent-space finrank consumed by the regularity criterion.

\begin{lemma}\leanok
[Stage 6.B (RHS) --- residue-tensor finrank from free rank]
  \label{lem:finrank_residue_kaehler_of_free_rank}
  \lean{AlgebraicGeometry.Scheme.finrank_residueField_tensor_kaehlerDifferential_of_free_rank_eq}
  Let \(R \to S_{\mathfrak q}\) be an algebra with \(S_{\mathfrak q}\) a nontrivial
  local ring whose K\"ahler differentials \(\Omega_{S_{\mathfrak q}/R}\) are
  \(S_{\mathfrak q}\)-free of rank \(n\). Then the residue-field base change
  \(\kappa \otimes_{S_{\mathfrak q}} \Omega_{S_{\mathfrak q}/R}\) has
  \(\kappa\)-dimension \(n\), where \(\kappa\) is the residue field of \(S_{\mathfrak q}\).
\end{lemma}

\begin{proof}
  Proved directly in Lean: base change preserves rank, and finite rank equals the
  natural number \(n\). A sub-step of
  \cref{lem:finrank_cotangentSpace_of_formallySmooth_residue}.
\end{proof}


\begin{lemma}[Stage 6.B --- cotangent isomorphism, maximal-ideal form]
  \label{lem:cotangent_iso_maximalIdeal_residue_kaehler}
  \lean{AlgebraicGeometry.Scheme.cotangent_iso_maximalIdeal_residue_tensor_kaehler_of_formallySmooth_residue}
  \uses{lem:cotangent_kahler_over_field}
  Under the closed-point hypotheses of \cref{lem:cotangent_kahler_over_field}
  (\(S_{\mathfrak q}\) formally smooth over \(R\), residue field \(\kappa\) formally
  smooth over \(R\), and \(\Omega_{\kappa/R} = 0\)), there is an
  \(S_{\mathfrak q}\)-linear isomorphism between the cotangent module of the maximal
  ideal \(\mathfrak m\) and the residue-field base change
  \(\kappa \otimes_{S_{\mathfrak q}} \Omega_{S_{\mathfrak q}/R}\). This restates the
  kernel-side cotangent isomorphism of \cref{lem:cotangent_kahler_over_field} with the
  canonical maximal-ideal cotangent space as its domain.
\end{lemma}

\begin{proof}
  Proved directly in Lean: transports the isomorphism of
  \cref{lem:cotangent_kahler_over_field} along the equality
  \(\ker(S_{\mathfrak q} \to \kappa) = \mathfrak m\).
\end{proof}

\begin{lemma}\leanok
[Stage 6.B --- cotangent-space finrank, formally-smooth residue form]
  \label{lem:finrank_cotangentSpace_of_formallySmooth_residue}
  \lean{AlgebraicGeometry.Scheme.finrank_cotangentSpace_of_formallySmooth_residue}
  \uses{lem:cotangent_iso_maximalIdeal_residue_kaehler, lem:finrank_residue_kaehler_of_free_rank}
  With the closed-point hypotheses of
  \cref{lem:cotangent_iso_maximalIdeal_residue_kaehler} and
  \(\Omega_{S_{\mathfrak q}/R}\) free of \(S_{\mathfrak q}\)-rank \(n\), the cotangent
  space \(\mathfrak m / \mathfrak m^2\) has \(\kappa\)-dimension \(n\).
\end{lemma}

\begin{proof}
  Proved directly in Lean: combine the \(\kappa\)-promoted cotangent isomorphism
  \cref{lem:cotangent_iso_maximalIdeal_residue_kaehler} with the residue-tensor
  finrank \cref{lem:finrank_residue_kaehler_of_free_rank}.
\end{proof}

\begin{lemma}\leanok
[Stage 6.B --- cotangent-space finrank at a $\bar k$-rational point]
  \label{lem:finrank_cotangentSpace_of_bijective_residue}
  \lean{AlgebraicGeometry.Scheme.finrank_cotangentSpace_of_bijective_algebraMap_residue}
  \uses{lem:finrank_cotangentSpace_of_formallySmooth_residue}
  Bundled closed-point form of
  \cref{lem:finrank_cotangentSpace_of_formallySmooth_residue}: if \(S_{\mathfrak q}\)
  is formally smooth over \(R\), \(\Omega_{S_{\mathfrak q}/R}\) is free of rank \(n\),
  and the structure map \(R \to \kappa\) to the residue field is bijective (as at a
  \(\bar k\)-rational closed point, by the Nullstellensatz), then
  \(\dim_\kappa \mathfrak m/\mathfrak m^2 = n\).
\end{lemma}

\begin{proof}
  Proved directly in Lean: bijectivity upgrades the residue field to formally smooth
  with vanishing differentials, reducing to
  \cref{lem:finrank_cotangentSpace_of_formallySmooth_residue}.
\end{proof}

\begin{lemma}
  [Stage 6.C --- Assembly: smooth standard-smooth stalk is regular local]
  \label{lem:stage6_regular_stalk_assembly}
  \uses{lem:rank_kaehler_localization_eq_relative_dim,
        lem:smooth_algebra_krull_dim_formula,
        lem:cotangent_kahler_over_field}
  Let \(\bar k\) be an algebraically closed field, let
  \(R \to S\) be a standard-smooth \(\bar k\)-algebra presentation of
  relative dimension \(n\) on an affine chart \(\Spec S \subseteq X\),
  and let \(z \in \Spec S\) be a point corresponding to a prime
  \(\mathfrak q \subset S\). Then the stalk \(S_{\mathfrak q}\) is a
  regular local ring of Krull dimension
  \(n - \mathrm{trdeg}_{\bar k}\,\kappa(\mathfrak q)\). In the
  closed-point case (\(\kappa(\mathfrak q) = \bar k\),
  \(\mathrm{trdeg} = 0\)) the dimension equals \(n\); in the
  codim-\(1\) generic-point case (\(\mathrm{trdeg} = n - 1\)) the
  dimension equals \(1\).
\end{lemma}

\begin{proof}
  Combine the three sub-lemmas:
  \begin{enumerate}
    \item By \cref{lem:rank_kaehler_localization_eq_relative_dim} (Stage~5b)
      the localised K\"ahler module \(\Omega_{S_{\mathfrak q} / R}\)
      is free of rank \(n\) over \(S_{\mathfrak q}\).
    \item By \cref{lem:smooth_algebra_krull_dim_formula} (Stage 6.A,
      Stacks 00OE) the Krull dimension of \(S_{\mathfrak q}\) equals
      \(\dim_z \Spec S - \mathrm{trdeg}_{\bar k}\,\kappa(\mathfrak q)\),
      which for a standard-smooth presentation of relative dimension
      \(n\) equals \(n - \mathrm{trdeg}_{\bar k}\,\kappa(\mathfrak q)\).
      (For the standard-smooth presentation the dimension
      \(\dim_z \Spec S = n\) follows from the global-complete-intersection
      structure: \(n\) variables modulo \(n - \dim\) equations, see
      Stacks \texttt{lemma-relative-global-complete-intersection-smooth}.)
    \item By \cref{lem:cotangent_kahler_over_field} (Stage 6.B,
      Stacks 02JK) the residue-field tensor \(\Omega_{S_{\mathfrak q}/R}
      \otimes_{S_{\mathfrak q}} \kappa(\mathfrak q)\) is isomorphic to
      \(\mathfrak m / \mathfrak m^2\) as \(\kappa(\mathfrak q)\)-vector
      spaces. Combined with (1), \(\dim_{\kappa(\mathfrak q)} \mathfrak
      m / \mathfrak m^2 = n\) at the closed-point case (and in general
      \(n - \mathrm{trdeg}_{\bar k}\,\kappa(\mathfrak q)\) by base-change
      of the standard-smooth presentation to the closed-point fibre).
    \item Compare (2) and (3): both \(\dim S_{\mathfrak q}\) and
      \(\dim_{\kappa(\mathfrak q)} \mathfrak m / \mathfrak m^2\) equal
      \(n - \mathrm{trdeg}_{\bar k}\,\kappa(\mathfrak q)\); the
      regularity criterion for Noetherian local rings (the equality
      \(\dim_\kappa \mathfrak m / \mathfrak m^2 = \dim R\) is the
      defining condition for a regular local ring, Stacks
      \texttt{definition-regular-local} \(=\) Mathlib
      \texttt{IsRegularLocalRing}) then gives \(S_{\mathfrak q}\)
      regular. The scheme-level transfer
      \(S_{\mathfrak q} = \mathcal O_{X, z}\) (Stage 3
      \texttt{IsLocalization.AtPrime} on the affine chart) completes
      the conclusion.
  \end{enumerate}
\end{proof}

\subsection{Stacks 00OE closed-point substrate (superseded decomposition)}
\label{subsec:stage6_iib_substrate_iter200}

An earlier decomposition of Stage 6.A
(\cref{lem:smooth_algebra_krull_dim_formula}, Stacks 00OE) reduced the
localised Krull-dimension formula at a closed point to a
Jacobian-regular-sequence witness for the relations of a submersive
presentation (Stacks 00SW / 00OW together with Matsumura Thm.~14.2). That
reduction is not needed: the closed-point case is proved unconditionally by the
Krull-height-theorem route ---
full height of maximal ideals of polynomial rings
(\cref{lem:mvPolynomial_height_eq_natCard_of_isMaximal}), the height
lower bound at maximal ideals of standard-smooth algebras
(\cref{lem:standard_smooth_le_height_of_isMaximal}), the localised
dimension lower bound
(\cref{lem:standard_smooth_le_ringKrullDim_localization}), and the
cotangent upper bound
(\cref{lem:regular_of_finrank_cotangent_le_dim}) --- which needs no
regular-sequence witness. The Matsumura regular-sequence lemma itself
survives as \cref{lem:matsumura_isRegular_of_linearIndependent_cotangent}.
The general (non-closed) point case does not require the
localisation theorem for regular local rings (Stacks \texttt{00OF}) at all:
regularity at an arbitrary prime is obtained directly, Serre-free, by the
trdeg--height route of \cref{subsec:smooth_prime_regularity}
(\cref{thm:standard_smooth_regular_at_prime}).

\subsection{Substrate lemmas for Stacks 00OE and the scheme bridges}
\label{subsec:codim1_substrate_lemmas}

The narrative above is realised on the Lean side by the named substrate lemmas
collected here. They feed the Stacks-00OE Krull-dimension formula
\cref{lem:smooth_algebra_krull_dim_formula} and the regular-stalk assembly.

\paragraph{Stacks 00OE localisation/height bridge.}

\begin{lemma}\leanok
[00OE Step 1 --- localised Krull dimension equals height]
  \label{lem:ringKrullDim_localization_eq_height_atPrime}
  \lean{AlgebraicGeometry.Scheme.ringKrullDim_localization_eq_height_atPrime}
  For a prime ideal \(\mathfrak m\) of a commutative ring \(R\) and any localisation
  \(A\) of \(R\) at \(\mathfrak m\), the Krull dimension of \(A\) equals the height of
  \(\mathfrak m\).
\end{lemma}

\begin{proof}
  Proved directly in Lean: re-export of Mathlib's
  \texttt{IsLocalization.AtPrime.ringKrullDim\_eq\_height}. Consumed by the
  closed-point dimension bound
  \cref{lem:standard_smooth_le_ringKrullDim_localization}.
\end{proof}

\paragraph{Closed-point form of the 00OE formula (Krull height theorem route).}

The following route proves the localised Krull-dimension formula at a maximal
ideal unconditionally: the dimension \emph{lower} bound comes from the full
height of the pulled-back maximal ideal in the polynomial ring via Krull's
height theorem, and the \emph{upper} bound is automatic from the cotangent
space. Combined with the
cotangent computation \cref{lem:finrank_cotangentSpace_of_bijective_residue},
this proves regularity of a standard-smooth algebra at every closed point.

\begin{lemma}\leanok
[Maximal ideals of polynomial rings over a field have full height]
  \label{lem:mvPolynomial_height_eq_natCard_of_isMaximal}
  \lean{MvPolynomial.height_eq_natCard_of_isMaximal}
  For a field \(k\) and a finite index set \(I\), every maximal ideal \(M\) of
  the polynomial ring \(k[X_i : i \in I]\) has height exactly \(\#I\).
\end{lemma}

\begin{proof}
  Induction on \(\#I\). For \(I = \emptyset\) the ring is the field \(k\), of
  Krull dimension \(0\), so every prime has height \(0\). For the induction
  step write \(k[X_i : i \in I \sqcup \{\ast\}] = R[X]\) with
  \(R = k[X_i : i \in I]\). Given a maximal ideal \(M \subseteq R[X]\), its
  contraction \(\mathfrak p = M \cap R\) is again maximal, because \(R\) is a
  finitely generated algebra over a field and hence a Jacobson ring. For a
  maximal ideal of \(R[X]\) over a Noetherian ring \(R\) one has
  \(\operatorname{ht} M = \operatorname{ht} \mathfrak p + 1\): the fibre
  \(R[X] \otimes_R \kappa(\mathfrak p) = \kappa(\mathfrak p)[X]\) is a
  principal ideal domain, so the image of \(M\) there has height \(1\), and
  going-down along the flat extension \(R \to R[X]\) splices the chains. The
  claim follows from the inductive hypothesis
  \(\operatorname{ht}\mathfrak p = \#I\).
\end{proof}

\begin{lemma}\leanok
[Height lower bound at maximal ideals of standard-smooth algebras]
  \label{lem:standard_smooth_le_height_of_isMaximal}
  \lean{Algebra.IsStandardSmoothOfRelativeDimension.natCast_le_height_of_isMaximal}
  Let \(k\) be a field and \(S\) a standard-smooth \(k\)-algebra of relative
  dimension \(n\). Then every maximal ideal \(\mathfrak m \subseteq S\) has
  height at least \(n\).
\end{lemma}

\begin{proof}
  \uses{lem:mvPolynomial_height_eq_natCard_of_isMaximal}
  Choose a submersive presentation \(S = P/I\) with
  \(P = k[X_1, \dots, X_N]\) and \(I = (f_1, \dots, f_c)\) generated by the
  \(c\) relations, so that \(n = N - c\); the submersive structure embeds the
  relation indices into the generator indices, so \(c \le N\) and
  \(n + c = N\). The preimage \(M \subseteq P\) of \(\mathfrak m\) is a maximal
  ideal containing \(I\), of height \(N\) by
  \cref{lem:mvPolynomial_height_eq_natCard_of_isMaximal}. Krull's height
  theorem, in the relative form
  \(\operatorname{ht} M \le \operatorname{ht}(M/I) + \mu(I)\) where
  \(\mu(I) \le c\) is the minimal number of generators of \(I\), gives
  \(N \le \operatorname{ht}(\mathfrak m) + c\) since \(M/I\) corresponds to
  \(\mathfrak m\) under \(P/I \simeq S\). Hence
  \(n \le \operatorname{ht}\mathfrak m\).
\end{proof}

\begin{lemma}\leanok
[Dimension lower bound at closed points of standard-smooth algebras]
  \label{lem:standard_smooth_le_ringKrullDim_localization}
  \lean{Algebra.IsStandardSmoothOfRelativeDimension.le_ringKrullDim_of_isLocalization_atPrime}
  Let \(k\) be a field, \(S\) a standard-smooth \(k\)-algebra of relative
  dimension \(n\), \(\mathfrak m \subseteq S\) a maximal ideal, and
  \(S_{\mathfrak m}\) a localisation of \(S\) at \(\mathfrak m\). Then
  \(n \le \dim S_{\mathfrak m}\).
\end{lemma}

\begin{proof}
  \uses{lem:standard_smooth_le_height_of_isMaximal,
        lem:ringKrullDim_localization_eq_height_atPrime}
  Combine \(\dim S_{\mathfrak m} = \operatorname{ht}\mathfrak m\)
  (\cref{lem:ringKrullDim_localization_eq_height_atPrime}) with the height
  lower bound \cref{lem:standard_smooth_le_height_of_isMaximal}.
\end{proof}

\begin{lemma}\leanok
[Regularity from the cotangent upper bound]
  \label{lem:regular_of_finrank_cotangent_le_dim}
  \lean{IsRegularLocalRing.of_finrank_cotangentSpace_le_ringKrullDim}
  Let \(A\) be a Noetherian local ring with maximal ideal \(\mathfrak m\) and
  residue field \(\kappa\). If
  \(\dim_\kappa \mathfrak m/\mathfrak m^2 \le \dim A\), then \(A\) is a
  regular local ring.
\end{lemma}

\begin{proof}
  By Nakayama's lemma the minimal number of generators \(\mu(\mathfrak m)\)
  equals \(\dim_\kappa \mathfrak m/\mathfrak m^2\), and Krull's height theorem
  gives \(\dim A \le \mu(\mathfrak m)\) in any Noetherian local ring. The
  hypothesis forces \(\mu(\mathfrak m) = \dim A\), which is the definition of
  regularity.
\end{proof}

\begin{lemma}\leanok
[Stacks 00TT at a rational closed point, algebra level]
  \label{lem:standard_smooth_regular_at_rational_point}
  \lean{AlgebraicGeometry.Scheme.isRegularLocalRing_localization_of_isStandardSmooth_of_bijective_residue}
  Let \(k\) be a field, \(S\) a standard-smooth \(k\)-algebra,
  \(\mathfrak m \subseteq S\) a maximal ideal, and \(S_{\mathfrak m}\) a
  localisation of \(S\) at \(\mathfrak m\) such that the composite
  \(k \to S_{\mathfrak m} \to \kappa(\mathfrak m)\) is bijective (a
  \(k\)-rational point). Then \(S_{\mathfrak m}\) is a regular local ring.
\end{lemma}

\begin{proof}
  \uses{lem:standard_smooth_le_ringKrullDim_localization,
        lem:regular_of_finrank_cotangent_le_dim,
        lem:finrank_cotangentSpace_of_bijective_residue,
        lem:module_free_kaehler_standardSmooth,
        lem:module_free_kaehler_localization,
        lem:rank_kaehler_localization_eq_relative_dim,
        lem:exists_standardSmoothOfRelativeDimension_of_standardSmooth}
  Let \(n\) be the relative dimension of \(S\) over \(k\)
  (\cref{lem:exists_standardSmoothOfRelativeDimension_of_standardSmooth}).
  The ring \(S\) is Noetherian (finite type over a field), hence so is
  \(S_{\mathfrak m}\), and \(S_{\mathfrak m}\) is formally smooth over \(k\)
  (a localisation of a smooth algebra). Its K\"ahler module
  \(\Omega_{S_{\mathfrak m}/k}\) is free of rank \(n\)
  (\cref{lem:module_free_kaehler_standardSmooth},
  \cref{lem:module_free_kaehler_localization},
  \cref{lem:rank_kaehler_localization_eq_relative_dim}), so by the conormal
  computation at the \(k\)-rational point
  (\cref{lem:finrank_cotangentSpace_of_bijective_residue})
  \(\dim_{\kappa} \mathfrak m/\mathfrak m^2 = n\). Since
  \(n \le \dim S_{\mathfrak m}\)
  (\cref{lem:standard_smooth_le_ringKrullDim_localization}), regularity
  follows from \cref{lem:regular_of_finrank_cotangent_le_dim}.
\end{proof}

\begin{lemma}\leanok
[Stacks 00TT at closed points over an algebraically closed field]
  \label{lem:standard_smooth_regular_at_closed_point_algclosed}
  \lean{AlgebraicGeometry.Scheme.isRegularLocalRing_localizationAtPrime_of_isStandardSmooth_of_isAlgClosed}
  Let \(\bar k\) be an algebraically closed field and \(S\) a standard-smooth
  \(\bar k\)-algebra. Then for every maximal ideal \(\mathfrak m \subseteq S\)
  the local ring \(S_{\mathfrak m}\) is regular.
\end{lemma}

\begin{proof}
  \uses{lem:standard_smooth_regular_at_rational_point}
  By Zariski's lemma the residue field \(S/\mathfrak m\) is a finite, hence
  algebraic, extension of \(\bar k\); as \(\bar k\) is algebraically closed,
  \(\bar k \to S/\mathfrak m\) is bijective, and composing with the
  isomorphism \(S/\mathfrak m \simeq \kappa(\mathfrak m)\) (valid since
  \(\mathfrak m\) is maximal) shows every maximal ideal is a
  \(\bar k\)-rational point. Conclude with
  \cref{lem:standard_smooth_regular_at_rational_point}.
\end{proof}

\begin{lemma}\leanok
[Step B.e --- standard-smooth algebras over an algebraically closed field are reduced]
  \label{lem:standard_smooth_reduced_algclosed}
  \lean{AlgebraicGeometry.Scheme.isReduced_of_isStandardSmooth_of_isAlgClosed}
  \uses{lem:standard_smooth_regular_at_closed_point_algclosed,
        lem:isDomain_of_regularLocal}
  Let \(\bar k\) be an algebraically closed field and \(S\) a standard-smooth
  \(\bar k\)-algebra. Then \(S\) is reduced.
\end{lemma}

\begin{proof}
  If \(S\) is the zero ring the claim is vacuous, so assume \(S \neq 0\).
  Reducedness is detected at localisations at maximal ideals: if \(x \in S\) is
  nilpotent and nonzero, its annihilator is a proper ideal, hence contained in
  some maximal ideal \(\mathfrak m\); the image of \(x\) in \(S_{\mathfrak m}\)
  is then nilpotent and nonzero, so \(S_{\mathfrak m}\) is not reduced. It
  therefore suffices to show every \(S_{\mathfrak m}\) is reduced. By
  \cref{lem:standard_smooth_regular_at_closed_point_algclosed} the local ring
  \(S_{\mathfrak m}\) is regular, by \cref{lem:isDomain_of_regularLocal} it is
  an integral domain, and integral domains are reduced.
\end{proof}

\begin{lemma}\leanok
[Step B.f --- a scheme smooth over an algebraically closed field is reduced]
  \label{lem:smooth_scheme_reduced_algclosed}
  \lean{AlgebraicGeometry.Scheme.isReduced_of_smooth_of_isAlgClosed}
  \uses{lem:exists_standardSmooth_at_of_smooth, lem:gammaSpecField_ringEquiv,
        lem:isLocalization_atPrime_stalk_of_affineOpen,
        lem:standard_smooth_reduced_algclosed}
  Let \(\bar k\) be an algebraically closed field and let \(X\) be a scheme with
  a smooth morphism \(X \to \Spec(\bar k)\). Then \(X\) is reduced.
\end{lemma}

\begin{proof}
  A scheme is reduced if and only if all of its stalks are reduced, so fix
  \(z \in X\). By \cref{lem:exists_standardSmooth_at_of_smooth} there are affine
  opens \(U \subseteq \Spec(\bar k)\) and \(V \ni z\) with \(V\) mapping into
  \(U\) such that \(\Gamma(\Spec \bar k, U) \to \Gamma(X, V)\) is standard
  smooth. Since \(U\) contains the image of \(z\) it is nonempty, so
  \cref{lem:gammaSpecField_ringEquiv} identifies
  \(\Gamma(\Spec \bar k, U) \simeq \bar k\); transporting the presentation along
  this isomorphism makes \(\Gamma(X, V)\) a standard-smooth
  \(\bar k\)-algebra, which is reduced by
  \cref{lem:standard_smooth_reduced_algclosed}. The stalk
  \(\mathcal O_{X,z}\) is a localisation of \(\Gamma(X, V)\) by
  \cref{lem:isLocalization_atPrime_stalk_of_affineOpen}, and a localisation of
  a reduced ring is reduced.
\end{proof}

\paragraph{Relative dimension and scheme-to-algebra bridges.}

\begin{lemma}\leanok
[Stage 6 sub-gap (i) --- relative dimension of a standard-smooth algebra]
  \label{lem:exists_standardSmoothOfRelativeDimension_of_standardSmooth}
  \lean{AlgebraicGeometry.Scheme.exists_isStandardSmoothOfRelativeDimension_of_isStandardSmooth}
  Every standard-smooth \(R\)-algebra \(S\) is standard smooth of some specific
  relative dimension \(n \in \mathbb N\).
\end{lemma}

\begin{proof}
  Proved directly in Lean: read the dimension off the underlying submersive
  presentation. A sub-step of \cref{lem:smooth_to_regular_local_ring}.
\end{proof}

\begin{lemma}\leanok
[Step B.a --- submersive presentation of a relative-dimension-$n$ algebra]
  \label{lem:exists_submersivePresentation_of_standardSmoothReldim}
  \lean{AlgebraicGeometry.Scheme.exists_submersivePresentation_of_isStandardSmoothOfRelativeDimension}
  \uses{lem:exists_standardSmoothOfRelativeDimension_of_standardSmooth}
  An \(R\)-algebra \(S\) that is standard smooth of relative dimension \(n\) admits an
  explicit submersive presentation \(P\) with \(P.\mathrm{dimension} = n\).
\end{lemma}

\begin{proof}
  Proved directly in Lean: re-export of the underlying-presentation field of the
  relative-dimension instance.
\end{proof}

\begin{lemma}\leanok
[Step B.b --- affine-open stalk is a localisation]
  \label{lem:isLocalization_atPrime_stalk_of_affineOpen}
  \lean{AlgebraicGeometry.Scheme.isLocalization_atPrime_stalk_of_affineOpen}
  For an affine open \(V \ni z\) of a scheme \(X\), the stalk \(\mathcal O_{X,z}\) is
  the localisation of the section ring \(\Gamma(X, V)\) at the prime ideal of \(z\).
\end{lemma}

\begin{proof}
  Proved directly in Lean: re-export of \texttt{IsAffineOpen.isLocalization\_stalk}.
\end{proof}

\begin{lemma}\leanok
[Step B.c helper --- nonempty open on a one-point space is everything]
  \label{lem:open_eq_top_of_subsingleton}
  \lean{AlgebraicGeometry.Scheme.open_eq_top_of_subsingleton}
  On a scheme whose underlying space is a subsingleton (at most one point), every
  nonempty open subset equals the whole space.
\end{lemma}

\begin{proof}
  Proved directly in Lean: the unique-point argument. A sub-step of
  \cref{lem:gammaSpecField_ringEquiv}.
\end{proof}

\begin{lemma}\leanok
[Step B.c --- sections of $\Spec \bar k$ over a nonempty open are $\bar k$]
  \label{lem:gammaSpecField_ringEquiv}
  \lean{AlgebraicGeometry.Scheme.gammaSpecField_ringEquiv}
  \uses{lem:open_eq_top_of_subsingleton}
  For a field \(\bar k\) and any nonempty open \(U\) of \(\Spec(\bar k)\), the section
  ring \(\Gamma(\Spec \bar k, U)\) is ring-isomorphic to \(\bar k\). This gives the
  algebraically-closed base ring of the standard-smooth presentation a clean
  \(\bar k\)-identification.
\end{lemma}

\begin{proof}
  Proved directly in Lean: \(\Spec(\bar k)\) has a single point, so
  \cref{lem:open_eq_top_of_subsingleton} forces \(U = \top\) and the sections collapse
  to the global sections \(\Gamma(\Spec \bar k, \top) \simeq \bar k\).
\end{proof}

\paragraph{Matsumura regular-sequence bridge (Stacks 00NQ).}

\begin{lemma}\leanok
[Step A1 bridge --- module quotient equals ring quotient]
  \label{lem:quotSMulTop_quotientRing_linearEquiv}
  \lean{AlgebraicGeometry.Scheme.quotSMulTop_quotientRing_linearEquiv}
  For \(r\) in a commutative ring \(A\), there is a canonical \((A/(r))\)-linear
  isomorphism between the module quotient \(A/(r \cdot A)\) and the ring quotient
  \(A/(r)\).
\end{lemma}

\begin{proof}
  Proved directly in Lean: compose the tensor description of the module quotient with
  the right-unitor and promote scalars along the surjective quotient map. A sub-step of
  \cref{lem:isRegular_cons_of_quotient_ring}.
\end{proof}

\begin{lemma}\leanok
[Step A1 bridge --- ring-quotient cons rule for regular sequences]
  \label{lem:isRegular_cons_of_quotient_ring}
  \lean{AlgebraicGeometry.Scheme.isRegular_cons_of_quotient_ring}
  \uses{lem:quotSMulTop_quotientRing_linearEquiv}
  Let \(r \in A\) be such that multiplication by \(r\) is injective on \(A\), and
  suppose the images of a list \(rs\) form a regular sequence over the ring quotient
  \(A/(r)\). Then \(r :: rs\) is a regular sequence over \(A\).
\end{lemma}

\begin{proof}
  Proved directly in Lean: the module-quotient cons rule transported across the
  identification \cref{lem:quotSMulTop_quotientRing_linearEquiv}. A sub-step of
  \cref{lem:matsumura_isRegular_of_linearIndependent_cotangent}.
\end{proof}

\begin{lemma}\leanok
[Step A1 --- cotangent independence descends to the quotient]
  \label{lem:matsumura_descent_cotangent}
  \lean{AlgebraicGeometry.Scheme.matsumura_descent_cotangent}
  Let \((A, \mathfrak m)\) be a local ring and \(f_0,\dots,f_m \in \mathfrak m\) with
  \(\kappa(A)\)-linearly independent cotangent classes. Then the tail
  \(f_1,\dots,f_m\), pushed into the quotient \(A' = A/(f_0)\), has
  \(\kappa(A')\)-linearly independent cotangent classes in
  \(\mathfrak m'/\mathfrak m'^2\).
\end{lemma}

\begin{proof}
  Proved directly in Lean: the cotangent map \(A \to A'\) is surjective with kernel
  spanned by the class of \(f_0\); a kernel-disjointness argument reduces a
  \(\kappa(A')\)-relation to a \(\kappa(A)\)-relation killed by the full independence
  hypothesis. A sub-step of
  \cref{lem:matsumura_isRegular_of_linearIndependent_cotangent}.
\end{proof}

\begin{lemma}\leanok
[Step A1 --- Matsumura: cotangent-independent elements form a regular sequence]
  \label{lem:matsumura_isRegular_of_linearIndependent_cotangent}
  \lean{AlgebraicGeometry.Scheme.matsumura_isRegular_of_linearIndependent_cotangent}
  \uses{lem:matsumura_descent_cotangent, lem:isRegular_cons_of_quotient_ring}
  Let \((A, \mathfrak m)\) be a Noetherian regular local ring and \(f_1,\dots,f_n \in
  \mathfrak m\) whose cotangent classes in \(\mathfrak m/\mathfrak m^2\) are
  \(\kappa(A)\)-linearly independent. Then \(f_1,\dots,f_n\) form a regular sequence in
  \(A\) (Matsumura, \emph{Commutative Ring Theory}, Thm.~14.2; Stacks 00NQ).
\end{lemma}

\begin{proof}
  Proved directly in Lean: induction on \(n\). The head \(f_1\) is a non-zero-divisor
  (a regular local ring is a domain), the cotangent independence descends to the
  quotient \(A/(f_1)\) by \cref{lem:matsumura_descent_cotangent}, and the
  ring-quotient cons rule \cref{lem:isRegular_cons_of_quotient_ring} reassembles the
  regular sequence.
\end{proof}

\subsection{Regularity at arbitrary primes (Serre-free Stacks 00TT)}
\label{subsec:smooth_prime_regularity}

The closed-point regularity of \cref{lem:standard_smooth_regular_at_closed_point_algclosed}
extends to \emph{every} prime of a standard-smooth algebra over a perfect
field, without invoking the localisation theorem for regular local rings
(Stacks tag 00OF, Serre's homological characterisation). The route combines a
transcendence-degree/height inequality in polynomial rings with the conormal
sequence at an arbitrary prime.

\begin{lemma}\leanok
[One-variable step of the trdeg--height inequality]
  \label{lem:polynomial_step_height_trdeg}
  \lean{Polynomial.step_height_trdeg_of_isPrime}
  % SOURCE: [Stacks Project], tag 00ON (dimension formula for going-down
  % extensions, Matsumura 13.B Theorem 19(2)), used through Mathlib's
  % Ideal.height_eq_height_add_of_liesOver_of_hasGoingDown.
  Let \(k\) be a field, \(R\) a Noetherian \(k\)-algebra, \(P \subseteq R[X]\)
  a prime ideal, and \(p = P \cap R\). Then either
  \[\operatorname{ht} p + 1 \le \operatorname{ht} P
    \quad\text{and}\quad
    \operatorname{trdeg}_k (R/p) \le \operatorname{trdeg}_k (R[X]/P),\]
  or
  \[\operatorname{ht} P = \operatorname{ht} p
    \quad\text{and}\quad
    \operatorname{trdeg}_k (R/p) + 1 \le \operatorname{trdeg}_k (R[X]/P).\]
\end{lemma}

\begin{proof}
  \leanok
  The extension \(R \to R[X]\) is free, hence flat, hence satisfies going-down,
  so the dimension formula (Stacks 00ON) gives
  \(\operatorname{ht} P = \operatorname{ht} p + \operatorname{ht} \bar P\) for
  the image \(\bar P\) of \(P\) in \(R[X]/pR[X] \cong (R/p)[X]\). Write
  \(A = R/p\) (a domain) and \(Q \subseteq A[X]\) for the prime corresponding
  to \(\bar P\); then \(Q \cap A = 0\) and \(A[X]/Q \cong R[X]/P\) as
  \(k\)-algebras. If \(Q = 0\), the heights agree and
  \(R[X]/P \cong A[X]\); superadditivity of transcendence degree in the tower
  \(k \to A \to A[X]\) together with \(\operatorname{trdeg}_A A[X] = 1\) gives
  \(\operatorname{trdeg}_k A[X] \ge \operatorname{trdeg}_k A + 1\). If
  \(Q \neq 0\), then \(0 \subsetneq Q\) is a chain of primes in the domain
  \(A[X]\), so \(\operatorname{ht} \bar P = \operatorname{ht} Q \ge 1\), while
  \(A\) still embeds into \(A[X]/Q\) because \(Q \cap A = 0\), whence
  \(\operatorname{trdeg}_k A \le \operatorname{trdeg}_k (A[X]/Q)\).
\end{proof}

\begin{lemma}\leanok
[trdeg--height inequality at primes of polynomial rings]
  \label{lem:mvpoly_trdeg_height}
  \lean{MvPolynomial.exists_le_trdeg_and_natCard_le_height_add}
  \uses{lem:polynomial_step_height_trdeg}
  Let \(k\) be a field, \(\iota\) a finite set, and
  \(P \subseteq k[x_i : i \in \iota]\) a prime ideal. Then there is
  \(d \in \mathbb N\) with
  \[d \le \operatorname{trdeg}_k \bigl(k[x_i]/P\bigr)
    \qquad\text{and}\qquad
    \#\iota \le \operatorname{ht} P + d.\]
  (Classically \(d\) may be taken to be the transcendence degree of the residue
  field and both inequalities are equalities; the witness form suffices for the
  regularity pipeline and requires no Noether normalisation.)
\end{lemma}

\begin{proof}
  \leanok
  Induction on \(\#\iota\). If \(\iota = \emptyset\) take \(d = 0\). For the
  step, write \(k[x_i : i \in \iota] = R[X]\) with
  \(R = k[x_i : i \in \iota\setminus\{i_0\}]\) and contract \(P\) to
  \(p = P \cap R\); the inductive hypothesis provides \(d\) with
  \(d \le \operatorname{trdeg}_k (R/p)\) and \(\#\iota - 1 \le
  \operatorname{ht} p + d\). In the height-jump case of
  \cref{lem:polynomial_step_height_trdeg},
  \(\#\iota \le (\operatorname{ht} p + 1) + d \le \operatorname{ht} P + d\)
  and the same \(d\) works; in the height-preserving case,
  \(\#\iota \le \operatorname{ht} P + (d+1)\) and
  \(d + 1 \le \operatorname{trdeg}_k (R[X]/P)\), so \(d+1\) works.
\end{proof}

\begin{lemma}\leanok
[trdeg--height inequality for standard-smooth algebras]
  \label{lem:standard_smooth_trdeg_height}
  \lean{Algebra.IsStandardSmoothOfRelativeDimension.exists_le_trdeg_and_natCast_le_height_add}
  \uses{lem:mvpoly_trdeg_height}
  Let \(S\) be a standard-smooth algebra of relative dimension \(n\) over a
  field \(k\) and \(q \subseteq S\) a prime ideal. Then there is
  \(d \in \mathbb N\) with \(d \le \operatorname{trdeg}_k (S/q)\) and
  \(n \le \operatorname{ht} q + d\).
\end{lemma}

\begin{proof}
  \leanok
  Choose a submersive presentation \(k[x_i : i \in \iota] \twoheadrightarrow S\)
  whose kernel is generated by \(\#\sigma\) relations, with
  \(n = \#\iota - \#\sigma\). Let \(M\) be the preimage of \(q\); then
  \(k[x_i]/M \cong S/q\), and \cref{lem:mvpoly_trdeg_height} provides \(d\)
  with \(d \le \operatorname{trdeg}_k (S/q)\) and
  \(\#\iota \le \operatorname{ht} M + d\). The relative form of Krull's height
  theorem bounds \(\operatorname{ht} M \le \operatorname{ht} q + \#\sigma\)
  (the same argument as in \cref{lem:standard_smooth_le_height_of_isMaximal});
  cancelling \(\#\sigma\) gives \(n \le \operatorname{ht} q + d\).
\end{proof}

\begin{lemma}\leanok
[K\"ahler rank equals transcendence degree over a perfect field]
  \label{lem:kaehler_rank_eq_trdeg}
  \lean{Algebra.rank_kaehlerDifferential_eq_trdeg}
  Let \(k\) be a perfect field and \(K\) a field extension of \(k\) that is
  essentially of finite type. Then
  \[\dim_K \Omega_{K/k} = \operatorname{trdeg}_k K.\]
\end{lemma}

\begin{proof}
  \leanok
  Since \(k\) is perfect, \(K\) admits a separating transcendence basis \(s\):
  writing \(F = k(s)\) for the rational-function subfield, the extension
  \(K/F\) is separable algebraic, hence formally \'etale, so base change along
  \(F \to K\) identifies \(\Omega_{K/k} \cong K \otimes_F \Omega_{F/k}\).
  Moreover \(F\) is the fraction field of the polynomial ring \(k[s]\), so
  \(\Omega_{F/k}\) is the localisation of \(\Omega_{k[s]/k}\), which is free
  with basis \(\{dx_i : i \in s\}\). Hence
  \(\dim_K \Omega_{K/k} = \#s = \operatorname{trdeg}_k K\).
\end{proof}

\begin{lemma}\leanok
[Conormal dimension identity at an arbitrary prime]
  \label{lem:conormal_finrank_identity}
  \lean{finrank_cotangentSpace_add_finrank_kaehler_residueField}
  % SOURCE: [Stacks Project], tag 02JK (the conormal sequence of a formally
  % smooth local algebra).
  Let \((S_{\mathfrak m}, \mathfrak m, \kappa)\) be a local \(R\)-algebra that
  is formally smooth over \(R\), whose residue field \(\kappa\) is also
  formally smooth over \(R\), and whose K\"ahler module
  \(\Omega_{S_{\mathfrak m}/R}\) is free of rank \(n\). Then
  \[\dim_\kappa \mathfrak m/\mathfrak m^2 + \dim_\kappa \Omega_{\kappa/R} = n.\]
\end{lemma}

\begin{proof}
  \leanok
  The conormal sequence
  \(\mathfrak m/\mathfrak m^2 \to \kappa \otimes_{S_{\mathfrak m}}
  \Omega_{S_{\mathfrak m}/R} \to \Omega_{\kappa/R} \to 0\)
  is exact, and formal smoothness of \(\kappa\) over \(R\) provides a
  retraction of its first map, which is therefore (split) injective. The
  middle term has \(\kappa\)-dimension \(n\) by freeness, and the last map is
  surjective; rank--nullity gives the identity.
\end{proof}

\begin{theorem}\leanok
[Standard-smooth algebras over a perfect field are regular at every prime]
  \label{thm:standard_smooth_regular_at_prime}
  \lean{isRegularLocalRing_of_isLocalization_atPrime_of_isStandardSmooth_of_perfectField}
  \uses{lem:standard_smooth_trdeg_height, lem:kaehler_rank_eq_trdeg,
        lem:conormal_finrank_identity, lem:regular_of_finrank_cotangent_le_dim,
        lem:ringKrullDim_localization_eq_height_atPrime,
        lem:module_free_kaehler_localization,
        lem:rank_kaehler_localization_eq_relative_dim}
  % SOURCE: [Stacks Project], tag 00TT (smooth algebras over fields), proved
  % here without tag 00OF (localisation of regular local rings).
  Let \(k\) be a perfect field, \(S\) a standard-smooth \(k\)-algebra of
  relative dimension \(n\), and \(q \subseteq S\) any prime ideal. Then any
  localisation \(S_q\) of \(S\) at \(q\) is a regular local ring.
\end{theorem}

\begin{proof}
  \leanok
  Write \(\kappa\) for the residue field of \(S_q\). The K\"ahler module
  \(\Omega_{S_q/k}\) is free of rank \(n\) (the localisation of the free
  module \(\Omega_{S/k}\), \cref{lem:module_free_kaehler_localization} and
  \cref{lem:rank_kaehler_localization_eq_relative_dim}); \(S_q\) is formally
  smooth over \(k\) (localisations of standard-smooth algebras are formally
  smooth), and \(\kappa\) is formally smooth over \(k\) since it is an
  essentially-finite-type field extension of a perfect field (separating
  transcendence bases exist). By \cref{lem:conormal_finrank_identity},
  \(\dim_\kappa \mathfrak m/\mathfrak m^2 = n - \dim_\kappa \Omega_{\kappa/k}\).
  By \cref{lem:kaehler_rank_eq_trdeg} and the embedding of \(S/q\) into
  \(\kappa\), \(\dim_\kappa \Omega_{\kappa/k} = \operatorname{trdeg}_k \kappa
  \ge \operatorname{trdeg}_k (S/q) \ge d\) for the witness \(d\) of
  \cref{lem:standard_smooth_trdeg_height}, which also gives
  \(n \le \operatorname{ht} q + d\). Hence
  \(\dim_\kappa \mathfrak m/\mathfrak m^2 \le n - d \le \operatorname{ht} q
  = \dim S_q\) (\cref{lem:ringKrullDim_localization_eq_height_atPrime}), and a
  Noetherian local ring whose cotangent dimension is at most its Krull
  dimension is regular (\cref{lem:regular_of_finrank_cotangent_le_dim}).
\end{proof}

\begin{lemma}\leanok
  [Smooth over \(\bar k\) \(\Rightarrow\) regular local stalk (Stacks 00TT)]
  \label{lem:smooth_to_regular_local_ring}
  \lean{AlgebraicGeometry.Scheme.isRegularLocalRing_stalk_of_smooth}
  \uses{lem:exists_standardSmooth_at_of_smooth, lem:gammaSpecField_ringEquiv,
        lem:isLocalization_atPrime_stalk_of_affineOpen,
        thm:standard_smooth_regular_at_prime}
  % SOURCE: [Stacks Project], tag 00TT --- Jacobian criterion
  % (smooth-implies-regular direction) (read from references/stacks-algebra.tex,
  % L38593--38611; tag verified against the Stacks website tag/00TT page,
  % "Lemma 10.140.3" in section 10.140 "Smooth algebras over fields").
  \textit{Source: [Stacks Project], tag 00TT.}
  Let \(X\) be a variety over an algebraically closed field \(\bar k\)
  (geometrically integral, separated, locally of finite type) and assume the
  structure morphism \(X \to \Spec(\bar k)\) is \emph{smooth}. Then for every
  point \(z \in X\) the stalk \(\mathcal O_{X, z}\) is a regular local ring.
\end{lemma}

\begin{proof}
  \leanok
  Choose an affine chart \(V \ni z\) on which the section ring
  \(\Gamma(X, V)\) is standard-smooth over
  \(\Gamma(\Spec \bar k, U) \cong \bar k\)
  (\cref{lem:exists_standardSmooth_at_of_smooth},
  \cref{lem:gammaSpecField_ringEquiv}). The stalk \(\mathcal O_{X,z}\) is the
  localisation of \(\Gamma(X, V)\) at the prime ideal of \(z\)
  (\cref{lem:isLocalization_atPrime_stalk_of_affineOpen}). Since \(\bar k\) is
  algebraically closed, it is perfect, and
  \cref{thm:standard_smooth_regular_at_prime} applies at every prime — closed
  or not — giving regularity of \(\mathcal O_{X,z}\).
\end{proof}

\begin{lemma}\leanok
[Smooth at a codim-1 point $\Rightarrow$ principal nonzero maximal ideal]
  \label{lem:smooth_codim_one_maximalIdeal_principal}
  \lean{AlgebraicGeometry.Scheme.smooth_codim_one_maximalIdeal_isPrincipal_and_ne_bot}
  \uses{lem:smooth_to_regular_local_ring, thm:ringKrullDim_stalk_eq_coheight}
  Let \(X\) be a smooth, geometrically-irreducible, separated, integral variety over an
  algebraically closed field \(\bar k\), and let \(z \in X\) be a point with
  \(\mathrm{coheight}\,z = 1\) (a codimension-one point). Then the maximal ideal of the
  stalk \(\mathcal O_{X,z}\) is principal and nonzero.
\end{lemma}

\begin{proof}
  Smoothness gives a regular local stalk via
  \cref{lem:smooth_to_regular_local_ring}; the coheight bridge
  \cref{thm:ringKrullDim_stalk_eq_coheight} gives Krull dimension \(1\); the
  cotangent-space characterisation of regularity then forces
  \(\dim_\kappa \mathfrak m/\mathfrak m^2 = 1\), i.e.\ a principal maximal ideal, which
  is nonzero since the stalk is not a field. This is the geometric core consumed by
  \cref{lem:smooth_codim_one_dvr}.
\end{proof}

\begin{lemma}


\leanok
  [Smooth \(\Rightarrow\) DVR at every codim-\(1\) point]
  \label{lem:smooth_codim_one_dvr}
  \lean{AlgebraicGeometry.Scheme.localRing_dvr_of_codim_one}
  \uses{lem:smooth_to_regular_local_ring, thm:ringKrullDim_stalk_eq_coheight,
        lem:smooth_codim_one_maximalIdeal_principal}
  % SOURCE: [Hartshorne], II.6, p.~130 (the definition ``regular in codimension one'',
  % together with the immediately preceding remark that on a nonsingular variety every
  % local ring is regular) (read from references/hartshorne-algebraic-geometry.pdf,
  % PDF page 147).
  % SOURCE QUOTE: "\textbf{Definition.} We say a scheme \(X\) is \emph{regular in
  % codimension one} (or sometimes \emph{nonsingular in codimension one}) if every
  % local ring \(\mathcal O_x\) of \(X\) of dimension one is regular."
  % SOURCE QUOTE: "The most important examples of such schemes are nonsingular
  % varieties over a field (I, \S5) and noetherian normal schemes. On a nonsingular
  % variety the local ring of every closed point is regular (I, 5.1), hence all the
  % local rings are regular, since they are localizations of the local rings of
  % closed points."
  % SOURCE QUOTE: "If \(Y\) is a prime divisor on \(X\), let \(\eta \in Y\) be its generic
  % point. Then the local ring \(\mathcal O_{\eta, X}\) of \(X\) at \(\eta\) is a discrete
  % valuation ring with quotient field \(K\), the function field of \(X\)."
  \textit{Source: Hartshorne, II.6, p.~130 (definition of regular in codimension one;
  the DVR consequence at a codim-1 point).}
  Let \(X\) be a nonsingular variety over \(\bar k\) and let \(\eta \in X\) be a point with
  \(\mathrm{codim}_X \overline{\{\eta\}} = 1\). Then the local ring \(\mathcal O_{X, \eta}\) is a
  discrete valuation ring with fraction field the function field \(K(X)\) of \(X\).
\end{lemma}

\begin{proof}

  By \cref{lem:smooth_to_regular_local_ring} (Stacks tag \texttt{00TT}, the
  Jacobian-criterion direction), the stalk \(\mathcal O_{X, x}\) is regular at
  every point \(x \in X\); in particular, for \(\eta \in X\) with
  \(\mathrm{codim}_X \overline{\{\eta\}} = 1\), the local ring
  \(\mathcal O_{X, \eta}\) is regular of Krull dimension \(1\). A regular local ring of Krull dimension \(1\) is exactly a
  discrete valuation ring (Atiyah--Macdonald, Proposition 9.2; or Stacks tag
  \texttt{00PD}). Its fraction field is the function field \(K(X)\) of \(X\) by
  irreducibility of \(X\) (the unique generic point of \(X\) lies above \(\eta\), and the
  stalk at the generic point of an integral scheme is the function field).

  \paragraph{Lean encoding scope.} The Lean target
  \texttt{AlgebraicGeometry.Scheme.localRing\_dvr\_of\_codim\_one} is the
  instance/lemma stating that \(\mathcal O_{X, \eta}\) is a Mathlib
  \texttt{IsDiscreteValuationRing} under the typeclass set
  \texttt{[Smooth\,(\dots)\,X.hom]\ [IsIntegral X.left]} and the hypothesis
  \texttt{Order.coheight\ \(\eta\) = 1}. The fraction-field identification with
  \texttt{X.functionField} is a downstream named instance
  \texttt{Scheme.fractionRing\_localRing\_eq\_functionField}. The regular-local-ring-of-dimension-1
  \(\Leftrightarrow\) DVR equivalence is shipped in Mathlib as
  \texttt{IsRegularLocalRing.isDiscreteValuationRing\_of\_dim\_one} (or analogous);
  the prover re-exports it rather than re-proving.
\end{proof}

\section{Milne Theorem 3.1: codimension-\(\geq 2\) extension}
\label{sec:milne_thm31}

The first input to Milne's Theorem~3.2 is the following: a rational map from a normal
variety to a \emph{complete} variety automatically extends across every codim-\(1\)
point.

\begin{lemma}\leanok
[Milne Theorem~3.1, codim-\(\geq 2\) conclusion]
  \label{thm:indeterminacy_codimGe2}
  \lean{AlgebraicGeometry.Scheme.RationalMap.indeterminacy_codimGe2_of_smooth_of_complete}
  \uses{def:indeterminacy_locus, lem:smooth_codim_one_dvr}
  % SOURCE: [Milne, Abelian Varieties], Theorem~3.1, \S I.3, p.~16 (the
  % valuative-criterion half of the proof).
  Let \(X\) be a nonsingular variety over \(\bar k\), \(Y\) a complete variety
  over \(\bar k\), and \(f \colon X \dashrightarrow Y\) a rational map of
  \(\bar k\)-varieties (i.e.\ defined over \(\bar k\): \(f\) followed by the
  structure morphism of \(Y\) is the structure morphism of \(X\), as rational
  maps). Then every point of the indeterminacy locus \(Z(f)\)
  (\cref{def:indeterminacy_locus}) has coheight at least \(2\); that is,
  \(Z(f)\) contains no point of codimension \(\le 1\).
\end{lemma}

\begin{proof}
  \leanok
  The generic point of \(X\) lies in \(\mathrm{Dom}(f)\), which is a dense
  open, so no point of coheight \(0\) lies in \(Z(f)\). Suppose \(z \in Z(f)\)
  had coheight \(1\). By \cref{lem:smooth_codim_one_dvr} the stalk
  \(\mathcal O_{X,z}\) is a discrete valuation ring with fraction field the
  function field \(K(X)\). The rational map \(f\) determines a morphism
  \(\Spec K(X) \to Y\), which together with the canonical morphism
  \(\Spec \mathcal O_{X,z} \to \Spec \bar k\) forms a commutative valuative
  square over the structure morphism of \(Y\) --- the square commutes
  because \(f\) is defined over \(\bar k\). The valuative criterion of
  properness for the complete variety \(Y\) produces a lift
  \(\Spec \mathcal O_{X,z} \to Y\). Spreading out (using that \(X\) is
  integral, hence germ-injective, and \(Y\) is of finite type over
  \(\bar k\)), \(f\) admits a representative defined on an open set
  containing \(z\), contradicting \(z \in Z(f)\).
\end{proof}


\begin{theorem}


\leanok
  [Milne Theorem~3.1: codim-\(\geq 2\) extension to a complete target]
  \label{thm:codim_one_extension}
  \lean{AlgebraicGeometry.Scheme.RationalMap.extend_of_codimOneFree_of_smooth}
  \uses{def:indeterminacy_locus, def:codim_one_indeterminacy,
        thm:indeterminacy_codimGe2, cor:regular_cohen_macaulay}
  % SOURCE: [Milne, Abelian Varieties], Theorem~3.1, \S I.3, p.~16
  % (read from references/abelian-varieties.pdf, PDF page 22).
  % SOURCE QUOTE: "\textsc{Theorem 3.1.} \emph{A rational map $\varphi \colon V
  % \dashrightarrow W\( from a normal variety \)V\( to a complete variety \)W$ is defined
  % on an open subset \(U\) of \(V\) whose complement \(V \smallsetminus U\) has codimension
  % \(\geq 2\).}"
  \textit{Source: Milne, \emph{Abelian Varieties}, Theorem~3.1, \S I.3, p.~16.}
  Let \(X\) be a nonsingular variety over \(\bar k\), let \(Y\) be a complete variety over
  \(\bar k\), and let \(f \colon X \dashrightarrow Y\) be a rational map of
  \(\bar k\)-varieties (defined over \(\bar k\)). Then the
  indeterminacy locus \(Z(f)\) (\cref{def:indeterminacy_locus}) has codimension
  \(\geq 2\) in \(X\); equivalently, \(f\) is codim-\(1\)-indeterminacy-free
  (\cref{def:codim_one_indeterminacy}).

  Specialising further: if in addition \(f\) \emph{is} codim-\(1\)-indeterminacy-free
  ``for free'' (i.e.\ no prime divisor of \(X\) is contained in \(Z(f)\)), then
  \(\mathrm{Dom}(f) = X\) and \(f\) extends to a regular morphism \(\widetilde f \colon X
  \to Y\). The extension is unique because \(X\) is reduced and \(\mathrm{Dom}(f)\) is dense
  in \(X\).
\end{theorem}

% SOURCE QUOTE PROOF: "\textsc{Proof.} Assume first that \(V\) is a curve. Thus we are
% given a nonsingular curve \(V\) and a regular map \(\varphi \colon U \to W\) from an
% open subset of \(V\) which we want to extend to \(V\). Consider the maps
% [diagram: \(U \hookrightarrow V\), $U \xrightarrow{u \mapsto (u, \varphi(u))} Z
% \subseteq V \times W\(, \)U \xrightarrow{\varphi} W$, with projections from
% \(V \times W\) to \(V\) and to \(W\)]
% Let \(U'\) be the image of \(U\) in \(V \times W\), and let \(Z\) be its closure. The image
% of \(Z\) in \(V\) is closed (because \(W\) is complete), and contains \(U\) (the composite
% \(U \to V\) is the given inclusion), and so \(Z\) maps onto \(V\). The maps $U \to U'
% \to U\( are isomorphisms. Therefore, \)Z \to V$ is a surjective map from a complete
% curve onto a nonsingular complete curve that is an isomorphism on open subsets.
% Such a map must be an isomorphism (complete nonsingular curves are determined by
% their function fields). The restriction of the projection map \(V \times W \to W\)
% to \(Z\ (\simeq V)\) is the extension of \(\varphi\) to \(V\) we are seeking.
%   The general case can be reduced to the one-dimensional case (using schemes). Let
% \(U\) be the largest subset on which \(\varphi\) is defined, and suppose that $V
% \smallsetminus U\( has codimension 1. Then there is a prime divisor \)Z\( in \)V
% \smallsetminus U\(. Because \)V$ is normal, its associated local ring is a discrete
% valuation ring \(\mathcal O_Z\) with field of fractions \(k(V)\). The map \(\varphi\)
% defines a morphism of schemes \(\operatorname{Spec}(k(V)) \to W\), which the
% valuative criterion of properness (Hartshorne 1977, II 4.7) shows extends to a
% morphism \(\operatorname{Spec}(\mathcal O_Z) \to W\). This implies that \(\varphi\)
% has a representative defined on an open subset that meets \(Z\) in a nonempty set,
% which is a contradiction."
\begin{proof}
  
  Following Milne's two-step proof.

  \emph{Step 1: rule out codim-\(1\) components.} Suppose for contradiction that \(Z(f)\)
  contains a prime divisor \(W \subseteq X\) (i.e.\ a codim-\(1\) irreducible component).
  Let \(\eta \in W\) be its generic point, so \(\mathrm{codim}_X \overline{\{\eta\}} = 1\). By
  \cref{lem:smooth_codim_one_dvr}, the local ring \(\mathcal O_{X, \eta}\) is a discrete
  valuation ring with fraction field the function field \(K(X)\). Any representative
  \((U, \varphi_U) \in f\) has \(U \cap \mathrm{Dom}(f) \neq \emptyset\) and so determines
  a \(\bar k\)-algebra map from the function field \(K(Y)\) of \(Y\) to \(K(X)\): explicitly,
  \(\varphi_U\) defines a morphism of schemes \(\operatorname{Spec}(K(X)) \to Y\) (since
  \(\operatorname{Spec}(K(X))\) is the generic point of \(X\)). By the valuative criterion
  of properness applied to the complete variety \(Y\)
  (\emph{cf.}\ Hartshorne~II.4.7, cited verbatim in Milne~3.1's proof), this morphism
  extends to a morphism \(\operatorname{Spec}(\mathcal O_{X, \eta}) \to Y\). Standard
  algebraic-geometry spreading then exhibits an open neighbourhood \(V\) of \(\eta\) in
  \(X\) on which \(f\) has a regular representative; in particular \(\eta \in
  \mathrm{Dom}(f)\), contradicting \(\eta \in W \subseteq Z(f)\). Hence no codim-\(1\)
  component of \(X\) lies in \(Z(f)\), i.e.\ \(Z(f)\) has codimension \(\geq 2\).

  \emph{Step 2: extend across codim-\(\geq 2\) points (the "codim-1-indeterminacy-free
  \(\Rightarrow\) extension" half).} Suppose now that \(Z(f)\) has codimension \(\geq 2\)
  but is non-empty (we want to deduce \(Z(f) = \emptyset\) under the additional input
  that ``\(f\) is already codim-1-indeterminacy-free, by hypothesis''). Pick a closed
  point \(x \in Z(f)\) with \(\mathrm{codim}_X \overline{\{x\}} \geq 2\). By
  \cref{cor:regular_cohen_macaulay} applied to the regular local ring
  \(\mathcal O_{X, x}\), we have
  \(\mathrm{depth}(\mathcal O_{X, x}) = \dim(\mathcal O_{X, x}) \geq 2\). The sheaf
  \(\mathcal O_X\) on the punctured neighbourhood \(V \smallsetminus \{x\}\)
  (any small affine neighbourhood \(V \ni x\)) has \(H^0_{\{x\}}(V, \mathcal O_X) = 0\)
  and \(H^1_{\{x\}}(V, \mathcal O_X) = 0\) (the depth-\(\geq 2\) characterisation of
  local-cohomology vanishing, Stacks tag \texttt{0AVF}; cf.\ Hartshorne~III.3.5 for
  the surface case). Pull back the coordinates of any affine chart of \(Y\) around the
  closure \(\overline{f(\mathrm{Dom}(f) \cap V)}\) along \(f\); the resulting sections of
  \(\mathcal O_X\) on \(V \cap \mathrm{Dom}(f)\) extend uniquely to sections on \(V\),
  yielding a regular representative of \(f\) on \(V\). Hence \(x \in \mathrm{Dom}(f)\),
  contradicting \(x \in Z(f)\). Therefore \(Z(f) = \emptyset\) and \(f\) extends to a
  regular morphism \(\widetilde f \colon X \to Y\).

  Uniqueness of \(\widetilde f\) is the standard reduced-and-separated agreement
  principle: two morphisms \(X \to Y\) that agree on the dense open \(\mathrm{Dom}(f)\)
  agree on all of \(X\).
\end{proof}

\begin{remark}
  \label{rmk:codim1_role_of_ab}
  \textbf{Where Auslander--Buchsbaum enters.} Step~1 of the proof uses only normality
  + the valuative criterion of properness (Hartshorne~II.4.7) and is characteristic-free.
  Step~2 uses the depth-\(\geq 2\) extension property at a codim-\(\geq 2\) regular
  point, which is exactly what \cref{cor:regular_cohen_macaulay} of
  \cref{chap:Albanese_AuslanderBuchsbaum} supplies (regular local of dimension \(d\)
  \(\Rightarrow\) Cohen--Macaulay of depth \(d\)). For our application the source \(X\) is a
  smooth variety, so the regular-local-ring hypothesis is automatic; the Cohen--Macaulay
  consequence is what feeds the local-cohomology vanishing.

  Concretely, when \(X\) is a smooth surface (e.g.\ \(\mathbb P^1 \times \mathbb P^1\) or
  \(C \times C\) for a smooth projective curve \(C\)), the codim-\(\geq 2\) closed points of
  \(X\) are exactly the closed points of \(X\) (because \(\dim X = 2\)), and at each such
  point \(\mathcal O_{X, x}\) is regular of dimension \(2\), hence Cohen--Macaulay of
  depth \(2\). The smooth-surface case is the one A.4.c will invoke when extending the
  symmetric-product rational map \(C^{(g)} \dashrightarrow A\) along the Abel--Jacobi
  birational equivalence to \(J = \Pic^0_{C/k}\).
\end{remark}

\section{Milne Lemma 3.3: indeterminacy locus of a map to a group variety is codim 1}
\label{sec:milne_lem33}

The second input to Milne's Theorem~3.2 is structural: when the target \(Y\) is a group
variety, the indeterminacy locus is either empty or of \emph{pure} codimension \(1\).
Combined with \cref{thm:codim_one_extension}'s ``codimension \(\geq 2\)'' conclusion,
this forces the locus to be empty.

\begin{lemma}


\leanok
  [Milne Lemma~3.3: pure-codim-\(1\) structure of the indeterminacy locus into a group
  variety]
  \label{lem:milne_codim1_indeterminacy}
  \lean{AlgebraicGeometry.Scheme.RationalMap.indeterminacy_pure_codim_one_into_grpScheme}
  \uses{def:indeterminacy_locus, def:codim1_cycles, def:order_at_point,
        def:principal_divisor}
  % SOURCE: [Milne, Abelian Varieties], Lemma~3.3, \S I.3, p.~17
  % (read from references/abelian-varieties.pdf, PDF page 23).
  % SOURCE QUOTE: "\textsc{Lemma 3.3.} \emph{Let \(\varphi \colon V \dashrightarrow G\)
  % be a rational map from a nonsingular variety to a group variety. Then either
  % \(\varphi\) is defined on all of \(V\) or the points where it is not defined form a
  % closed subset of pure codimension 1 in \(V\) (i.e., a finite union of prime
  % divisors).}"
  \textit{Source: Milne, \emph{Abelian Varieties}, Lemma~3.3, \S I.3, p.~17.}
  Let \(X\) be a nonsingular variety over \(\bar k\), let \(G\) be a group variety over
  \(\bar k\) (a smooth group scheme of finite type), and let \(f \colon X \dashrightarrow
  G\) be a rational map of \(\bar k\)-varieties (defined over \(\bar k\)).
  Then either \(f\) is defined on all of \(X\), or the indeterminacy
  locus \(Z(f) \subseteq X\) is a closed subset of pure codimension \(1\) (i.e.\ a finite
  union of prime divisors).

  Equivalently: it is \emph{never} the case that \(f\) has a non-empty indeterminacy
  locus all of whose irreducible components have codimension \(\geq 2\).
\end{lemma}

% SOURCE QUOTE PROOF: "\textsc{Proof.} Define a rational map
%   \(\Phi \colon V \times V \dashrightarrow G\), $\;(x, y) \mapsto \varphi(x) \cdot
%   \varphi(y)^{-1}$.
% More precisely, if \((U, \varphi_U)\) represents \(\varphi\), then \(\Phi\) is the rational
% map represented by
%   $U \times U \xrightarrow{\varphi_U \times \varphi_U} G \times G
%   \xrightarrow{\mathrm{id} \times \mathrm{inv}} G
%   \xrightarrow{m} G$.
% Clearly \(\Phi\) is defined at a diagonal point \((x, x)\) if \(\varphi\) is defined at
% \(x\), and then \(\Phi(x, x) = e\). Conversely, if \(\Phi\) is defined at \((x, x)\), then
% it is defined on an open neighbourhood of \((x, x)\); in particular, there will be an
% open subset \(U\) of \(V\) such that \(\Phi\) is defined on \(\{x\} \times U\). After
% possible replacing \(U\) by a smaller open subset (not necessarily containing \(x\)),
% \(\varphi\) will be defined on \(U\). For \(u \in U\), the formula
%   \(\varphi(x) = \Phi(x, u) \cdot \varphi(u)\)
% defines \(\varphi\) at \(x\). Thus \(\varphi\) is defined at \(x\) if and only if \(\Phi\) is
% defined at \((x, x)\). The rational map \(\Phi\) defines a map
%   \(\Phi^\ast \colon \mathcal O_{G, e} \to k(V \times V)\).
% Since \(\Phi\) sends \((x, x)\) to \(e\) if it is defined there, it follows that \(\Phi\)
% is defined at \((x, x)\) if and only if
%   \(\operatorname{Im}(\mathcal O_{G, e}) \subset \mathcal O_{V \times V, (x, x)}\).
% Now \(V \times V\) is nonsingular, and so we have a good theory of divisors (AG,
% Chapter 12). For a nonzero rational function \(f\) on \(V \times V\), write
%   \(\operatorname{div}(f) = \operatorname{div}(f)_0 - \operatorname{div}(f)_\infty\),
% with \(\operatorname{div}(f)_0\) and \(\operatorname{div}(f)_\infty\) effective divisors
% --- note that \(\operatorname{div}(f)_\infty = \operatorname{div}(f^{-1})_0\). Then
%   $\mathcal O_{V \times V, (x, x)} = \{f \in k(V \times V) \mid
%   \operatorname{div}(f)_\infty \text{ does not contain } (x, x)\} \cup \{0\}$.
% Suppose \(\phi\) is not defined at \(x\). Then for some $f \in \operatorname{Im}
% (\varphi^\ast)\(, \)(x, x) \in \operatorname{div}(f)_\infty\(, and clearly \)\Phi$ is
% not defined at the points \((y, y) \in \Delta \cap \operatorname{div}(f)_\infty\).
% This is a subset of pure codimension one in \(\Delta\) (AG 9.2), and when we identify
% it with a subset of \(V\), it is a subset of \(V\) of codimension one passing through
% \(x\) on which \(\varphi\) is not defined."
\begin{proof}
  
  Following Milne's proof of Lemma~3.3. The argument has four sub-steps. Throughout we
  fix a representative \((U, \varphi_U)\) of \(f\).

  \emph{Sub-step 1: the difference rational map.} Define
  \[
    \Phi \colon X \times X \dashrightarrow G,
    \qquad
    \Phi(x, y) \;:=\; \varphi(x) \cdot \varphi(y)^{-1},
  \]
  via the composition
  \[
    U \times U
    \;\xrightarrow{\varphi_U \times \varphi_U}\; G \times G
    \;\xrightarrow{\mathrm{id} \times \mathrm{inv}}\; G \times G
    \;\xrightarrow{m}\; G,
  \]
  where \(\mathrm{inv} \colon G \to G\) is the group-variety inverse and
  \(m \colon G \times G \to G\) is the group law. The map \(\Phi\) is rational by
  composition: its domain is at least \(U \times U \subseteq X \times X\), which is dense
  because \(U\) is dense in \(X\).

  \emph{Sub-step 2: \(\Phi\)-defined at \((x, x)\) \(\Leftrightarrow\) \(\varphi\)-defined at
  \(x\).} If \(\varphi\) is defined at \(x\), the explicit formula \(\Phi(x, x) = e\) (the
  identity of \(G\)) shows \(\Phi\) is defined at \((x, x)\). Conversely, if \(\Phi\) is
  defined at \((x, x)\), openness of \(\mathrm{Dom}(\Phi)\) in \(X \times X\) produces an
  open set \(W \ni (x, x)\) on which \(\Phi\) is defined; by projection (taking the slice
  \(\{x\} \times W' \subseteq W\) for some open \(W' \subseteq X\), possibly shrinking
  \(W'\) so \(\varphi\) is defined on \(W'\) as well --- \(W'\) need not contain \(x\)), the
  identity
  \[
    \varphi(x) \;=\; \Phi(x, u) \cdot \varphi(u),
    \qquad u \in W',
  \]
  exhibits a regular representative of \(\varphi\) at \(x\). Hence \(x \in
  \mathrm{Dom}(\varphi)\) if and only if \((x, x) \in \mathrm{Dom}(\Phi)\).

  \emph{Sub-step 3: rephrase via pullback of local rings.} The rational map \(\Phi\)
  pulls \(\mathcal O_{G, e}\) back to a \(\bar k\)-subalgebra of the function field
  \(K(X \times X)\):
  \[
    \Phi^\ast \colon \mathcal O_{G, e} \;\longrightarrow\; K(X \times X).
  \]
  Because \(\Phi\) sends \((x, x)\) to \(e\) wherever it is defined, \(\Phi\) is defined at
  \((x, x)\) if and only if \(\Phi^\ast(\mathcal O_{G, e}) \subseteq \mathcal O_{X
  \times X, (x, x)}\) (i.e.\ every local function on \(G\) near \(e\) has no pole at
  \((x, x)\) after pullback).

  \emph{Sub-step 4: pole locus is pure codimension \(1\).} The variety \(X \times X\) is
  nonsingular (smoothness is stable under products over \(\bar k\)), so the Weil-divisor
  apparatus of \cref{def:codim1_cycles} and the order map of
  \cref{def:order_at_point} apply on \(X \times X\). For any nonzero rational function
  \(f \in K(X \times X)^\ast\) decompose its divisor as
  \[
    \mathrm{div}(f) \;=\; \mathrm{div}(f)_0 - \mathrm{div}(f)_\infty,
  \]
  with \(\mathrm{div}(f)_0\) and \(\mathrm{div}(f)_\infty\) effective and
  \(\mathrm{div}(f)_\infty = \mathrm{div}(f^{-1})_0\) (Milne's notation, matching the
  project's \(\div(f) = \sum \ord_Y(f) \cdot Y\) of
  \cref{def:principal_divisor}: \(\mathrm{div}(f)_0\) collects the prime divisors with
  \(\ord_Y(f) > 0\) and \(\mathrm{div}(f)_\infty\) those with \(\ord_Y(f) < 0\)). The local
  ring at \((x, x)\) can be characterised as
  \[
    \mathcal O_{X \times X, (x, x)}
    \;=\; \bigl\{\, f \in K(X \times X) \;\big|\;
    \mathrm{div}(f)_\infty \text{ does not contain } (x, x) \,\bigr\} \cup \{0\}.
  \]
  Suppose \(f \in K(X \times X)^\ast\) is not in \(\mathcal O_{X \times X, (x, x)}\);
  equivalently, \((x, x) \in \mathrm{supp}(\mathrm{div}(f)_\infty)\). The pole divisor
  \(\mathrm{div}(f)_\infty\) is a finite sum of prime divisors of \(X \times X\) (i.e.\ a
  closed subset of pure codimension \(1\) in \(X \times X\)). Its intersection with the
  diagonal \(\Delta_X \subseteq X \times X\) is a closed subset of pure codimension \(1\)
  in \(\Delta_X\) (this is Milne's reference to AG~9.2, the standard fact that a prime
  divisor of \(X \times X\) either contains \(\Delta_X\) or meets it in pure codimension
  \(1\); on a nonsingular variety this is an instance of Krull's principal-ideal theorem
  applied to the local equation of \(\Delta_X\) inside \(X \times X\)).

  \emph{Conclusion.} Suppose \(f \colon X \dashrightarrow G\) is not defined on all of
  \(X\). Pick \(x \in Z(f)\). By Sub-steps~2--3, there is some \(f_0 \in
  \Phi^\ast(\mathcal O_{G, e})\) with \((x, x) \in \mathrm{supp}(\mathrm{div}(f_0)_\infty)\).
  By Sub-step~4 the intersection \(\mathrm{supp}(\mathrm{div}(f_0)_\infty) \cap
  \Delta_X\) is pure codimension \(1\) in \(\Delta_X\), passing through \((x, x)\).
  Identifying \(\Delta_X\) with \(X\) via the diagonal isomorphism, the corresponding
  closed subset of \(X\) has pure codimension \(1\) and passes through \(x\). By Sub-step~2,
  every point of this codim-\(1\) subset is a point where \(\Phi\) is undefined on the
  diagonal, hence where \(\varphi\) is undefined. Thus \(Z(\varphi)\) contains a pure
  codim-\(1\) closed subset passing through \(x\); since \(x \in Z(\varphi)\) was arbitrary,
  every irreducible component of \(Z(\varphi)\) has codimension exactly \(1\), i.e.\
  \(Z(\varphi)\) is of pure codimension \(1\).
\end{proof}

\begin{remark}
  \label{rmk:milne_lem33_charfree}
  \textbf{Characteristic-free.} The above proof uses only: (i) the group law on \(G\)
  and its inverse (formal smoothness of \(G\) as a group object); (ii) the Weil-divisor
  apparatus on the nonsingular variety \(X \times X\) (Hartshorne~II.6, or
  \cref{chap:RiemannRoch_WeilDivisor} in the project's notation); (iii) the
  decomposition of a divisor on a nonsingular variety into effective pieces, and
  Krull's principal-ideal theorem on \(X \times X\). None of these has a characteristic
  restriction, so \cref{lem:milne_codim1_indeterminacy} holds over every
  algebraically closed base field, in arbitrary characteristic. (Milne states the
  lemma in this generality at \S I.3, p.~17.)
\end{remark}

\section{Weil-divisor obstruction to codim-1 extension}
\label{sec:weil_obstruction}

The proof of \cref{thm:codim_one_extension}'s Step~1 (codim-\(1\) ruling-out) goes
through the valuative criterion. In the project's Weil-divisor language
(\cref{chap:RiemannRoch_WeilDivisor}) the same obstruction has a cleaner phrasing:
the extension fails at a prime divisor \(W\) if and only if some affine coordinate of
\(f\) around the image \(\overline{f(\eta_W)}\) has a strictly negative order
\(\ord_W(\cdot) < 0\). This is the form the Lean prover needs when consuming
\cref{def:order_at_point} and \cref{def:principal_divisor}.


\begin{remark}
  \label{rmk:weil_obstruction_application}
  \textbf{Why this form is the prover-facing one.} In the consumer chapter
  \cref{chap:Albanese_Thm32RationalMapExtension} the Lean prover combines
  \cref{thm:codim_one_extension} with \cref{lem:milne_codim1_indeterminacy}; in
  isolation, the abstract codim-\(\geq 2\) statement
  \cref{thm:codim_one_extension} suffices and \texttt{thm:weil_divisor_obstruction}
  is not invoked. The Weil-divisor form is the bridge to the project's
  \(\mathrm{Scheme.WeilDivisor}\) API, which is the form A.4.a's own internal sub-proofs
  use when establishing \emph{specific} extensions in genus-\(0\) classifying-map
  arguments and in the symmetric-product Albanese chapter
  \cref{chap:Albanese_AlbaneseUP}.
\end{remark}


\begin{remark}
  \label{rmk:mem_domain_partial_map_reshuffle_scope}
  \textbf{Relationship to \texttt{thm:weil_divisor_obstruction}.}
  \texttt{lem:mem_domain_partial_map_reshuffle} records the definitional half of
  the Hartshorne-II.6 ``regular = no pole'' reformulation that
  \texttt{thm:weil_divisor_obstruction} states in its full \(\ord_W \geq 0\)
  form: it provides a partial-map representative through a prime-divisor generic
  point. The full \(\ord_W \geq 0\) form additionally requires the pullback of a
  rational map to the function field, which is the content of
  \texttt{thm:weil_divisor_obstruction}.
\end{remark}

\section{Lean encoding}
\label{sec:codim1_lean_encoding}

The chapter's declarations correspond, in order, to the following Lean targets:

\begin{enumerate}
  \item \texttt{AlgebraicGeometry.Scheme.RationalMap.indeterminacyLocus} ---
    \cref{def:indeterminacy_locus}, a method on the Mathlib
    \texttt{Scheme.RationalMap} structure. Implementation: \texttt{X.carrier
    \textbackslash{} f.domain} as a \texttt{Set X}, paired with an
    \texttt{IsClosed} witness (either bundled or as a separate
    \texttt{Scheme.RationalMap.isClosed\_indeterminacyLocus} lemma).
  \item \texttt{AlgebraicGeometry.Scheme.RationalMap.CodimOneFree} ---
    \cref{def:codim_one_indeterminacy}, a \texttt{Prop}-valued predicate. Encoding:
    \texttt{\(\forall\) \(\eta\) : X, Order.coheight \(\eta\) = 1 \(\to\) \(\eta\) \(\in\)
    f.domain} (matching the \texttt{Order.coheight} idiom of
    \cref{def:prime_divisor}).
  \item \texttt{AlgebraicGeometry.Scheme.localRing\_dvr\_of\_codim\_one} ---
    \cref{lem:smooth_codim_one_dvr}, a lemma (or instance) producing
    \texttt{IsDiscreteValuationRing (Scheme.LocalRing.at\ X \(\eta\))} under the
    typeclass set \texttt{[Smooth (\(\Spec \bar k\)) X] [IsIntegral X]} and the
    hypothesis \texttt{Order.coheight \(\eta\) = 1}.
  \item \texttt{AlgebraicGeometry.Scheme.RationalMap.extend\_of\_codimOneFree\_of\_smooth}
    --- \cref{thm:codim_one_extension}, the existence-of-extension theorem. Signature
    (informal): given a smooth variety \(X\), a complete variety \(Y\), and
    \(f\)~codim-\(1\)-indeterminacy-free, produce a unique regular morphism
    \(\widetilde f \colon X \to Y\) extending \(f\). Consumes
    \cref{cor:regular_cohen_macaulay} from
    \cref{chap:Albanese_AuslanderBuchsbaum} as a Mathlib-style import.
  \item \texttt{AlgebraicGeometry.Scheme.RationalMap.indeterminacy\_pure\_codim\_one\_into\_grpScheme}
    --- \cref{lem:milne_codim1_indeterminacy}. Signature (informal): given a smooth
    variety \(X\), a group scheme \(G\) over \(\bar k\) that is a variety (smooth, finite
    type, separated), and \(f \colon X \dashrightarrow G\) rational, conclude that
    \(Z(f)\) is either empty or of pure codimension \(1\) in \(X\).
  \item \texttt{AlgebraicGeometry.Scheme.RationalMap.mem\_domain\_iff\_exists\_partialMap\_through\_point}
    --- the definitional reshuffle pinned at
    \texttt{lem:mem_domain_partial_map_reshuffle}. The substantive
    Hartshorne-II.6 ``regular = no pole'' Weil-divisor reformulation lives
    at \texttt{thm:weil_divisor_obstruction}. Re-uses
    \texttt{AlgebraicGeometry.Scheme.RationalMap.order} from
    \cref{chap:RiemannRoch_WeilDivisor} (the order map at a prime divisor).
\end{enumerate}


\section{Out of scope}
\label{sec:codim1_out_of_scope}

This chapter packages \emph{only} the two extension inputs and the Weil-divisor
reformulation needed for A.4.a. The following adjacent content is explicitly out of
scope:
\begin{itemize}
  \item \textbf{The Auslander--Buchsbaum formula and the regular \(\Rightarrow\)
    Cohen--Macaulay corollary.} Proved in \cref{chap:Albanese_AuslanderBuchsbaum}
    as \cref{thm:auslander_buchsbaum} and \cref{cor:regular_cohen_macaulay}; this
    chapter only consumes them (in Step~2 of \cref{thm:codim_one_extension}; see
    \cref{rmk:codim1_role_of_ab}).
  \item \textbf{Milne Theorem~3.2 itself.} The combination of
    \cref{thm:codim_one_extension} and \cref{lem:milne_codim1_indeterminacy} into the
    statement ``a rational map from a nonsingular variety to an abelian variety
    extends'' lives in \cref{chap:Albanese_Thm32RationalMapExtension} as
    \cref{thm:rational_map_to_av_extends}; this chapter only supplies the two inputs.
  \item \textbf{Milne Theorem~3.4 / Corollary~3.7 / Theorem~3.8 / Proposition~3.9 /
    Proposition~3.10.} The remainder of Milne~\S I.3 (the ``map between AVs is
    translated homomorphism'' family of results, and the
    \(\mathbb A^1 \dashrightarrow A\) / unirational\(\dashrightarrow A\) constancy
    statements) is downstream of \cref{thm:rational_map_to_av_extends} but is not
    needed by A.4 in the project's current routing: the genus-\(0\) base case is
    handled via the proven Rigidity Lemma + \(\mathbb G_m\)-scaling shortcut
    (\texttt{chap:RigidityKbar}), and the positive-genus arm consumes only the bare
    Theorem~3.2 statement.
  \item \textbf{Weil-divisor degree on \(X \times X\) for \(X\) a curve.} The proof of
    \cref{lem:milne_codim1_indeterminacy} invokes the Weil-divisor apparatus on the
    smooth surface \(X \times X\) but does \emph{not} need the degree map of
    \texttt{def:divisor_degree} (which is curve-specific); only the codim-\(1\)
    decomposition \(\div(f) = \div(f)_0 - \div(f)_\infty\) on a nonsingular surface is
    used. The bilinear degree formula on \(X \times X\) for \(X\) a curve (intersection
    numbers, Hartshorne~V.1) is out of scope.
  \item \textbf{Cartier divisors and the Cartier/Weil comparison.} On a smooth variety
    they coincide and the chapter uses Weil throughout; the Cartier formalism is not
    introduced.
  \item \textbf{Generalisation to non-algebraically-closed base fields.} The
    project's positive-genus arm runs over \(\bar k\) throughout; descent of the
    extension to a non-algebraically-closed base, if ever needed, would require
    Galois-descent machinery and is not in scope.
  \item \textbf{Resolution of indeterminacy by blow-ups} (Hironaka-style, for higher
    indeterminacy codimensions on a smooth variety with a non-group-variety target):
    out of scope. The combination of \cref{thm:codim_one_extension} and
    \cref{lem:milne_codim1_indeterminacy} for \(G\) a group variety is sharper than any
    blow-up-based approach because it produces an extension on the original \(X\)
    rather than on a blow-up; no blow-up apparatus is needed.
\end{itemize}
